% This file is generated from the corresponding Obsidian note.
% Source: Teoricos/Teoremas Corazon.md
\chapter{Teoremas corazón}
\label{ch:teoremas-corazon}
\section{Vectores tangentes}
\begin{bookproposition}{Localidad de los vectores tangentes}
Sea $M$ una variedad suave, $p\in M$ y $v\in T_pM$. Si $f,g\in C^\infty(M)$ coinciden en un abierto de $p$, entonces

\[
v(f)=v(g).
\]
\end{bookproposition}
\begin{bookproof}{}
\begin{enumerate}
\item Es equivalente demostrar que si $f-g=h\in C^\infty(M)$ se anula en un abierto $U$ de $p$, entonces $v(h)=0$.
\item Se toma una función campana $\beta$ con $\beta\equiv 1$ en un abierto $V\subseteq U$ de $p$ y soporte contenido en $U$.
\item Entonces $\beta h\equiv 0$ en todo $M$ por que si $x\in U$ $h(x)=0$ y si $x\in U^{c}$ entonces $\beta(x)=0$ entonces
\end{enumerate}
\[
v(\beta h)=0.
\]
por linealidad de $v$ 4. Pero por Leibniz,

\[
0=v(\beta h)=\beta(p)\,v(h)+h(p)\,v(\beta)=v(h)
\]
ya que $\beta(p)=1$ y $h(p)=0$. 5. Luego $v(h)=0$.

\bookqed
\end{bookproof}
\begin{bookproposition}{Independencia lineal de los vectores coordenados}
El conjunto de vectores coordenados

\[
\left\{\left.\frac{\partial}{\partial x_1}\right|_p,\dots,\left.\frac{\partial}{\partial x_n}\right|_p\right\}
\]
es linealmente independiente en $T_pM$.

\end{bookproposition}
\begin{bookproof}{}
6. Las funciones coordenadas $x_i$ están definidas solo en $U$, ósea están en $C^{\infty}(U)$ 7. Pero $\left.\frac{\partial}{\partial x_n}\right|_p$ esta en $C^{\infty}(M)$ así que primero se extienden a funciones suaves globales $\widetilde{x}_i\in C^\infty(M)$ que coinciden con $x_i$ en un abierto $V$ de $p$ usando \hyperref[blk:405b16]{Extensión de funciones suaves} (obtenemos múltiples abiertos y tomamos la intersección no vacía por que esta $p$) 8. Ahora notamos que

\[
\left.\frac{\partial}{\partial x_i}\right|_{p}\widetilde{x}_{j}=\left.\frac{\partial}{\partial x_i}\right|_{p}x_{j}
\]
porque $\widetilde{x}_{j}$ y $x_j$ coinciden en $V$ un abierto de $p$, por \hyperref[blk:b713ca]{Localidad de los vectores tangentes}. Luego aplicando la definición del vector coordenado,

\[
\left.\frac{\partial}{\partial x_i}\right|_{p}x_{j}=\left.\frac{\partial}{\partial y_i}\right|_{\varphi(p)}(x_{j}\circ\varphi^{-1}).
\]
\begin{enumerate}[start=9]
\item Y recordemos la notacion $(U,\varphi=(x_{1},\ldots,x_{n}))$ osea que $x_{j}\circ\varphi ^{-1}(q_{1},\ldots,q_{n})=x_{j}(p)=q_{j}$ otra manera de decir que $x_{j}\circ\varphi ^{-1}=\pi_{j}$ la proyeccion
\item Y es directo ver que
\end{enumerate}
\[
\left.\frac{\partial}{\partial y_i}\right|_{\varphi(p)}\pi_{j}=\delta_{ij}
\]
\begin{enumerate}[start=11]
\item Entonces tenemos
\end{enumerate}
\[
t_1\left.\frac{\partial}{\partial x_1}\right|_p+\cdots+t_n\left.\frac{\partial}{\partial x_n}\right|_p=0
\]
aplicando $\widetilde{x}_j$ a esta combinación lineal se obtiene

\[
0=t_1\delta_{1j}+\cdots+t_n\delta_{nj}=t_j
\]
y por lo tanto todos los coeficientes son cero. 12. Mostrando la independencia lineal.

\bookqed
\end{bookproof}
\section{El espacio tangente de una carta}
\begin{bookproposition}{Regla de la cadena}
Si $F:M\to N$ y $G:N\to P$ son suaves, entonces

\[
d(G\circ F)_p=(dG)_{F(p)}\circ dF_p.
\]
\end{bookproposition}
\begin{bookproof}{}
\begin{enumerate}
\item Sea $v\in T_pM$ y $h\in C^\infty(P)$. Entonces como la composición es asociativa
\end{enumerate}
\[
\bigl(d(G\circ F)_p(v)\bigr)(h)=v(h\circ (G\circ F))=v(h\circ G\circ F)
\]
\begin{enumerate}[start=2]
\item Por otro lado, mirando $dF_{p}(v)$ como un vector tangente y aplicandolo a $h\circ G$
\end{enumerate}
\[
\bigl(((dG)_{F(p)}\circ dF_p)(v)\bigr)(h)=((dG)_{F(p)}(dF_{p}(v)))(h)=(dF_p(v))(h\circ G)=v((h\circ G)\circ F)=v(h\circ G\circ F).
\]
\begin{enumerate}[start=3]
\item Como ambas derivaciones actúan igual sobre toda función $h\in C^\infty(P)$, son iguales.
\item Y luego como ambas actuan sobre todo $v\in T_{p}M$ son iguales
\end{enumerate}
\bookqed
\end{bookproof}
\section{Funciones de corte y funciones campana}
\begin{bookcorollary}{Función que vale 1 sobre un cerrado}
Sea $U$ un abierto de $M$ y $A$ un cerrado de $M$ con

\[
A\subseteq U
\]
Entonces existe una función suave

\[
f:M\to\mathbb R
\]
tal que

\[
0\le f\le 1,\qquad f|_A\equiv 1,\qquad \operatorname{supp}f\subseteq U
\]
\end{bookcorollary}
\begin{bookproof}{}
\begin{enumerate}
\item Como $A$ es cerrado. Considerar al cubrimiento de $M$ dado por
\end{enumerate}
\[
\{U,M\setminus A\}
\]
\begin{enumerate}[start=2]
\item Por \hyperref[blk:eafe2b]{Existencia de particiones de la unidad} existe una particion de la unidad $\{\rho,\sigma\}$ subordinada a $\{U,M\setminus A\}$ con
\end{enumerate}
\[
\operatorname{supp}\rho\subseteq U,\qquad \operatorname{supp}\sigma\subseteq M\setminus A
\]
\begin{enumerate}[start=3]
\item Para $p\in A$ se tiene $\sigma(p)=0$ porque $\operatorname{supp}\sigma\subseteq M\setminus A$. Usando la parte 3 de la \hyperref[blk:f5db1b]{definición de partición de la unidad}, tenemos $\rho(p)+\sigma(p)=1$, luego
\end{enumerate}
\[
\rho(p)=1
\]
\begin{enumerate}[start=4]
\item Tomamos $f=\rho$ que es la funcion que buscabamos
\end{enumerate}
\bookqed
\end{bookproof}
\begin{bookproposition}{Extensión de funciones suaves definidas en cerrados}
Sea $M$ una variedad suave, $A\subseteq M$ cerrado y

\[
F:A\to\mathbb R^k
\]
una función suave en el sentido anterior. Entonces, para todo abierto $U$ con $A\subseteq U$ existe una función suave

\[
\widetilde F:M\to\mathbb R^k
\]
funcion suave tal que

\[
\widetilde F|_A=F,
\qquad
\operatorname{supp}\widetilde F\subseteq U.
\]
\end{bookproposition}
\begin{bookproof}{}
\begin{enumerate}
\item Por \hyperref[blk:cf02ec]{Función suave sobre un subconjunto arbitrario}, para cada $p\in A$ elegimos un abierto $W_p\subseteq U$ y una extensión suave local
\end{enumerate}
\[
\widetilde F_p:W_p\to\mathbb R^k
\]
que coincide con $F$ en $W_{p}\cap A$ 2. Sin perdida de generalidades podemos suponer $W_{p}\subseteq U$ (En otro caso, cambiamos $W_{p}$ por el abierto $W_{p}\cap U$)\\
3. Ahora tenemos el conjunto

\[
\{W_p\}_{p\in A}\cup\{M\setminus A\}
\]
que es un cubrimiento abierto de $M$. 4. Por \hyperref[blk:eafe2b]{Existencia de particiones de la unidad} existe una partición de la unidad

\[
\{\rho_p\}_{p\in A}\cup\{\rho_0\}
\]
subordinada estrictamente a este cubrimiento, con

\[
\operatorname{supp}\rho_p\subseteq W_p, \qquad \operatorname{supp}\rho_0\subseteq M\setminus A
\]
\begin{enumerate}[start=5]
\item Como hicimos en el ejemplo anterior, definimos
\end{enumerate}
\[
\widetilde F=\sum_{p\in A}\rho_p\,\widetilde F_p
\]
donde estamos abusando de notación y entendemos el producto

\[
(\rho_p\widetilde{F}_p)(q)=\begin{cases}\rho_p(q)\,\widetilde{F}_p(q)&\text{si }q\in W_p,\\0&\text{si }q\notin W_p.\end{cases}
\]
\begin{enumerate}[start=6]
\item Veamos que cada término $\rho_p\widetilde F_p$ extendido por $0$ fuera de $W_p$ es suave sobre todo $M$. Si $q\in W_p$, la suavidad es trivial porque allí es producto de funciones suaves. Si $q\notin W_p$, entonces como $\operatorname{supp}\rho_p\subseteq W_p$, también $q\notin\operatorname{supp}\rho_p$. Como $\operatorname{supp}\rho_p$ es cerrado, $M\setminus\operatorname{supp}\rho_p$ es un abierto que contiene a $q$, y en ese abierto la extensión vale $0$; por tanto también es suave alrededor de $q$.
\item Tenemos que $\widetilde{F}$ es suave: como $\{\operatorname{supp}\rho_p\}$ es localmente finita, alrededor de cada punto de $M$ la suma es finita; y por el paso 6 cada término de esa suma finita es una función suave.
\item Ahora, si $q\in A$,
\end{enumerate}
\[
\widetilde{F}(q)=\sum_{p\in A}\rho_p(q)\widetilde{F}_p(q)=\sum_{p\in A}\rho_p(q)F(q)
\]
porque $\widetilde{F}_p$ coincide con $F$ en $W_p\cap A$, 9. Y seguimos

\[
\sum_{p\in A}\rho_p(q)F(q)=\left(\sum_{p\in A}\rho_p(q)\right)F(q)=\left(1-\rho_0(q)\right)F(q)
\]
donde el segundo igual vale por que $\sum_{p\in A}\rho_{p}+\rho_{0}=1$ 10. Y como $\operatorname{supp}\rho_0\subseteq M\setminus A$, tenemos

\[
\left(1-\rho_0(q)\right)F(q)=F(q).
\]
\begin{enumerate}[start=11]
\item Por último veamos $\operatorname{supp}\widetilde{F}\subseteq U$: Sea $q\in\operatorname{supp}\widetilde{F}$. Como $\{\operatorname{supp}\rho_p\}\cup\{\operatorname{supp}\rho_0\}$ es localmente finita, existe $W_q$ abierto de $M$ con $q$ que corta un número finito de miembros de $\{\operatorname{supp}\rho_p\}\cup\{\operatorname{supp}\rho_0\}$, digamos $\operatorname{supp}\rho_{p_1},\ldots,\operatorname{supp}\rho_{p_k}$ y $\operatorname{supp}\rho_0$.
\item Esto implica que debe haber un $j$ tal que $q\in\operatorname{supp}\rho_{p_j}$, pues en otro caso, existirían abiertos $W_1,\ldots,W_k$ de $q$ donde $\rho_{p_1},\ldots,\rho_{p_k}$ se anulan, y así
\end{enumerate}
\[
W_0\cap W_1\cap\cdots\cap W_k
\]
es un abierto de $q$ donde $\widetilde{F}$ vale cero que es absurdo 13. Como $q\in\operatorname{supp}\rho_{p_j}$ y $\operatorname{supp}\rho_{p_j}\subseteq W_{p_j}\subseteq U$, entonces $q\in U$. Por lo tanto, $\operatorname{supp}\widetilde F\subseteq U$.

\bookqed
\end{bookproof}
\section{Repaso de algebra lineal}
\begin{bookproposition}{Funciones suaves e independientes forman carta}
Sea $M$ una variedad suave de dimensión $m$. Sea $U$ un abierto de $M$ y $p\in U$. Si $\{y_1,\ldots,y_m\}$ es un conjunto de funciones suaves sobre $U$ \hyperref[blk:conjunto-funciones-independientes]{independientes en p} entonces existe un abierto $V$ de $M$ en $p$, con $V\subseteq U$, tal que

\[
(V,\psi=(y_1,\ldots,y_m))
\]
es una carta en $p$.

\end{bookproposition}
\begin{bookproof}{}
\begin{enumerate}
\item Se define
\end{enumerate}
\[
\psi:U \subseteq M \to \mathbb R^m,\qquad q \mapsto (y_1(q),\ldots,y_m(q)).
\]
La cual resulta suave. 2. Queremos ver que $(d\psi)_p$ es un isomorfismo. 3. Como $\dim M=m$ y $\{(dy_1)_p,\ldots,(dy_m)_p\}$ es linealmente independiente, entonces

\[
\mathcal B^*=\{(dy_1)_p,\ldots,(dy_m)_p\}
\]
es base de $T_p^*M$. 4. Sea

\[
\mathcal B=\{v_1,\ldots,v_m\}
\]
la respectiva base dual de $B^{*}$ en $T_pM$. 5. Usamos $\mathcal B$ para calcular $(d\psi)_p$.

\[
\begin{aligned}(d\psi)_p(v_j)&=\sum_{i=1}^m a_{ij}\frac{\partial}{\partial r_i}\Big|_{\psi(p)}\\&=\sum_{i=1}^m (d\psi)_{p}(v_{j})(r_{i})\frac{\partial}{\partial r_i}\Big|_{\psi(p)}\\&=\sum_{i=1}^m (v_{j})(r_{i}\circ\psi)\frac{\partial}{\partial r_i}\Big|_{\psi(p)}\\&=\sum_{i=1}^m (v_{j})(y_{i})\frac{\partial}{\partial r_i}\Big|_{\psi(p)}\\&=\sum_{i=1}^m (dy_{i})_{p}(v_{j})\frac{\partial}{\partial r_i}\Big|_{\psi(p)}\\&=\frac{\partial}{\partial r_i}\Big|_{\psi(p)}\end{aligned}
\]
\begin{enumerate}[start=6]
\item Luego, $(d\psi)_p$ manda base en base, por lo tanto es un isomorfismo y por el teorema de la funcion inversa existe un abierto $V$ de $p$ tal que $\psi(V)$ es abierto en $\mathbb{R}^{n}$ y $\psi|_{V}:V\rightarrow\psi(V)$ es un difeo
\end{enumerate}
\bookqed
\end{bookproof}
\section{Lema de factorización MUY IMPORTANTE}
\begin{bookcorollary}{}
Sean $M$ y $N$ variedades suaves de dimensión $m$ y $n$, respectivamente. Sea

\[
F:M\to N
\]
suave y supongamos que

\[
(dF)_p:T_pM\to T_{F(p)}N
\]
es inyectiva, y por tanto $m\le n$. (Osea $F$ es inmersion) Si $(V,\varphi=(x_1,\ldots,x_n))$ es una carta suave de $N$ en $F(p)$, entonces existen

\[
i_1,\ldots,i_m\in\{1,\ldots,n\}
\]
tales que

\[
\psi=(x_{i_1}\circ F,\ldots,x_{i_m}\circ F)
\]
define una carta suave de $M$ en cierto abierto de $p$.

\end{bookcorollary}
\begin{bookproof}{}
\begin{enumerate}
\item Para usar el \hyperref[blk:49431b]{Teórico 11}, alcanza con ver que
\end{enumerate}
\[
\{d(x_1\circ F)_p,\ldots,d(x_n\circ F)_p\}
\]
genera $T_p^\ast M$. 2. Como

\[
(dF)_p:T_pM\to T_{F(p)}N
\]
es inyectiva, entonces

\[
(dF)^\ast_{F(p)}:T_{F(p)}^\ast N\to T_p^\ast M
\]
es sobreyectiva por \hyperref[blk:dual-inyectiva-sobreyectiva]{dual de una transformacion inyectiva}. 3. Como $(V,\varphi=(x_1,\ldots,x_n))$ es una carta suave de $N$ en $F(p)$ sucede que

\[
\left\{(dx_1)_{F(p)},\ldots,(dx_n)_{F(p)}\right\}
\]
es base de $T_{F(p)}^\ast N$ por \hyperref[blk:base-dual-coordenadas]{base dual de las coordenadas}. Entonces sus imágenes por $(dF)^\ast_{F(p)}$ generan $T_p^\ast M$ por el paso 2. 4. Pero

\[
(dF)^\ast_{F(p)}\big((dx_i)_{F(p)}\big)=d(x_i\circ F)_p.
\]
Luego los $d(x_i\circ F)_p$ generan $T_p^\ast M$, como queríamos.

\bookqed
\end{bookproof}
\begin{bookproposition}{Lema de factorizacion}
Sea $F:M\rightarrow N$ suave. Y sea $(P,g)$ subvariedad de $N$ tal que $F(M)\subseteq g(P)$. Ademas sea $\widehat{F}:M\rightarrow P$ la funcion $\widehat{F}(x)=y$ si $F(x)=g(y)$ la cual es la unica funcion de $M$ en $P$ que hace conmutar el diagrama 
\bookimage{assets/pasted-image-20260429160319.png}{Pasted image 20260429160319}
 Entonces si $\widehat{F}$ es continua es suave

\end{bookproposition}
\begin{bookproof}{}
\begin{enumerate}
\item Sea $p\in M$ y veamos que $\widehat F$ es suave en un entorno de $p$.
\item Sea $q=\widehat F(p)$, es decir, $q$ es el único elemento de $P$ tal que
\end{enumerate}
\[
F(p)=g(q).
\]
\begin{enumerate}[start=3]
\item Sea $(V,\psi=(y_1,\ldots,y_n))$ una carta suave de $N$ alrededor de $F(p)$.
\item Como $g:P\to N$ es subvariedad, tenemos
\end{enumerate}
\[
(dg)_q:T_qP\to T_{g(q)}N
\]
es inyectiva, y por \hyperref[blk:c4d494]{Corolario}, existe un entorno $W$ de $q$ e índices $i_1,\ldots,i_k\in\{1,\ldots,n\}$, con $k=\dim P$, tal que

\[
(W,\varphi=(y_{i_1}\circ g,\ldots,y_{i_k}\circ g))
\]
es carta suave de $P$. 5. Sea

\[
\pi:\mathbb R^n\to\mathbb R^k,\qquad (s_1,\ldots,s_n)\mapsto (s_{i_1},\ldots,s_{i_k}),
\]
la cual es suave. 6. La ventaja de introducir esta función es que podemos expresar la carta $(W,\varphi)$ como

\[
(W,\pi\circ\psi\circ g).
\]
\begin{enumerate}[start=7]
\item Como $\widehat F$ es continua, tenemos
\end{enumerate}
\[
U=\widehat F^{-1}(W)
\]
es un abierto de $M$ que contiene a $p$. 8. Para ver que $\widehat{F}$ es suave en $p$ alcanza ver que $\widehat{F}|_{U}$ es suave y para esto componemos con la carta $\varphi$ (basta componer solo a un lado por que si vemos eso luego dada carta $\psi$ tendremos que $\varphi\circ\widehat{F}|_{U}\circ\psi$ es composicion de suaves)

\[
\varphi\circ\widehat{F}|_{U}=\pi\circ\psi\circ g\circ\widehat{F}|_{U}=\pi\circ\psi\circ F|_{U}
\]
esto ultimo vale por que $g(\widehat{F}|_{U}(x))=g(y)=F(x)$ por definicion 9. $\pi$ es suave por definicion $\psi\circ F|_{U}$ es suave por que $F$ es suave entonces $\varphi\circ\widehat{F}|_{U}$ resulta suave. Luego $\widehat{F}$ es suave

\bookqed
\end{bookproof}
\begin{bookproposition}{Unicidad de subvariedad via inclusion}
Sea $N$ variedad suave y sea $A\subseteq N$. Supongamos que $A$ tiene una topología $\tau$ no necesariamente la topología relativa. Entonces existe a lo sumo una estructura de variedad diferenciable sobre $A$ compatible con $\tau$ tal que

\[
(A,i)
\]
es subvariedad, donde

\[
i:A\to N
\]
es la inclusión. (De hecho podria no haber ninguna)

\end{bookproposition}
\begin{bookproof}{}
10. Sean $\mathcal F_1,\mathcal F_2$ dos estructuras de variedad diferenciable sobre $A$ compatibles con la topología $\tau$. 11. Para ver que son iguales, alcanza con ver que

\[
\operatorname{Id}:(A,\mathcal F_1)\to (A,\mathcal F_2)
\]
\[
\operatorname{Id}:(A,\mathcal F_2)\to (A,\mathcal F_1)
\]
es difeomorfismo. 12. Usemos el lema de factorizacion 
\bookimage{assets/pasted-image-20260429215527.png}{Pasted image 20260429215527}
 13. Como $F=i$, es claro que la (unica) $\widehat{F}$ que cumple \hyperref[blk:1258d0]{Lema de factorizacion} es

\[
\widehat F=\operatorname{Id}:(A,\mathcal F_1)\to(A,\mathcal F_2)
\]
\begin{enumerate}[start=14]
\item Como la topología que tienen $(A,\mathcal F_1)$ y $(A,\mathcal F_2)$ es la misma, tenemos que $\widehat F=\operatorname{Id}$ es continua, y por tanto $\widehat F$ es suave.
\item Cambiando los roles, también tenemos que $\operatorname{Id}:(A,\mathcal F_2)\to(A,\mathcal F_1)$ es suave.
\item Por tanto, $\operatorname{Id}:(A,\mathcal F_1)\to(A,\mathcal F_2)$ es un difeomorfismo, y entonces $\mathcal F_1=\mathcal F_2$
\item Se necesita que sea difeomorfismo para comparar cartas de ambas estructuras: si $(U,\varphi)\in\mathcal F_1$ y $(V,\psi)\in\mathcal F_2$, entonces en la intersección el cambio de coordenadas se puede escribir usando la identidad como
\end{enumerate}
\[
\psi\circ\varphi^{-1}=\psi\circ\operatorname{Id}\circ\varphi^{-1}.
\]
Como $\operatorname{Id}:(A,\mathcal F_1)\to(A,\mathcal F_2)$ es suave, esta composición es suave. Cambiando los roles se obtiene la suavidad de $\varphi\circ\psi^{-1}$, y por eso las cartas son compatibles.

\bookqed
\end{bookproof}
\section{Extension de funciones suaves}
\begin{bookproposition}{Extension local}
Sea $(M,i)$ una inmersion en $N$ y sea $p\in M$ y sea $g:M\to\mathbb R$ funcion suave. Entonces, existen un abierto $U$ de $M$ en $p$, un abierto $V$ de $N$ en $i(p)$ con $i(U)\subseteq V$ (cuidado, $M$ no necesariamente tiene topologia heredada) y una funcion

\[
\widetilde g:V\subseteq N\to\mathbb R
\]
suave que extiende a $g$, es decir, $\widetilde g\circ i=g$ para todo $q\in U$. 
\bookimage{assets/pasted-image-20260505230207.png}{Pasted image 20260505230207}


\end{bookproposition}
\begin{bookproof}{}
\begin{enumerate}
\item Sea $m=\dim M$ y $n=\dim N$ si $m=n$ sale gratis (por que?)
\item Como $(M,i)$ es una inmersion, sabemos por la formula local de una inmersion que existen cartas $(U,\varphi)$ y $(V,\psi)$ centradas en $p$ e $i(p)$ tales que
\end{enumerate}
\[
\psi\circ i\circ\varphi^{-1}(x_1,\ldots,x_m)=(x_1,\ldots,x_m,0,\ldots,0)
\]
\begin{enumerate}[start=3]
\item Para poder proyectar adecuadamente, tenemos $\epsilon$ tal que $C_\epsilon^m(0)\subseteq\varphi(U)$ y $C_\epsilon^n(0)\subseteq\psi(V)$. Luego llamamos
\end{enumerate}
\[
U=\varphi^{-1}(C_\epsilon^m(0))\quad \text{y}\quad V=\psi^{-1}(C_\epsilon^n(0))
\]
\begin{enumerate}[start=4]
\item Como hemos visto en otras pruebas, tenemos $i(U)\subseteq V$ pues aplicando $\psi$ de ambos lados, esto equivale a
\end{enumerate}
\[
\psi\circ i\circ\varphi^{-1}(C_\epsilon^m(0))\subseteq C_\epsilon^n(0)
\]
y esto ultimo es cierto. 5. Si $(x_1,\ldots,x_m)\in C_\epsilon^m(0)$, entonces $|x_i|<\epsilon$, por lo tanto $(x_1,\ldots,x_m,0,\ldots,0)\in C_\epsilon^n(0)$. 6. Sea $\pi:\mathbb R^n\to\mathbb R^m$, $(y_1,\ldots,y_m,y_{m+1},\ldots,y_n)\mapsto(y_1,\ldots,y_m)$. 7. Nos animamos a proponer $\widetilde g:V\subseteq N\to\mathbb R$ dada por

\[
\widetilde g=g\circ\varphi^{-1}\circ\pi\circ\psi
\]
\begin{enumerate}[start=8]
\item Como
\end{enumerate}
\[
q\in V\xrightarrow{\psi}\psi(q)\in C_{\epsilon}^{n}(0) \xrightarrow{\pi}\pi(\psi(q))\in C_{\epsilon}^{m}(0) \xrightarrow{\varphi^{-1}}\varphi^{-1}(\pi(\psi(q)))\in U\xrightarrow{g}g\circ\varphi^{-1}\circ\pi\circ\psi(q)=\tilde{g}(q)
\]
, la funcion $\widetilde g$ claramente es funcion suave (por composicion de suaves) sobre $V$ . 9. Ademas si $q\in U$,

\[
\widetilde g\circ i(q)=\tilde g\circ i\circ \varphi\circ\varphi ^{-1}(q)=g\circ\varphi^{-1}\circ\pi\circ\psi\circ i\circ\varphi ^{-1}\circ\varphi(q)=g\circ\varphi ^{-1}\circ Id_{C_{\epsilon}^{m} (0)}\circ\varphi(q)=g(q)
\]
\bookqed
\end{bookproof}
\begin{bookproposition}{Caracterizaciones de suavidad}
Sea $X:M \to TM$ un campo vectorial. Son equivalentes:

\begin{enumerate}
\item $X$ es suave.
\item Para toda carta suave $(U,\varphi=(x_1,\ldots,x_n))$ de $M$, las funciones $a_i:U \to \mathbb R$ dadas por
\end{enumerate}
\[
X_p=\sum_{i=1}^n a_i(p)\frac{\partial}{\partial x_i}|_p
\]
son suaves. 3. Para todo abierto $V$ de $M$ y toda funcion $f \in C^\infty(V)$, la funcion $Xf:V \to \mathbb R$, $(Xf)(q)=X_q f$, es suave. 4. Para toda funcion $f \in C^\infty(M)$, se cumple $Xf \in C^\infty(M)$.

\end{bookproposition}
\begin{bookproof}{}
\begin{itemize}
\item $1 \Rightarrow 2$.
\end{itemize}
\begin{enumerate}
\item Tenemos la carta $(U,\varphi=(x_{1},\ldots,x_{n}))$ la cual induce $(\widetilde{U}=\pi ^{-1}(U),\widetilde{\varphi})$ dada por
\end{enumerate}
\[
\widetilde\varphi(p,v)=(x_1(p),\ldots,x_n(p),v_1,\ldots,v_n)
\]
donde $v_i$ son las coordenadas del vector $v$ en la base $\{ \frac{\partial}{\partial x_{1}}|_{p},\ldots,\frac{\partial}{\partial x_{n}}|_{p} \}$\\
2. Como $X$ es suave,

\[
\widetilde\varphi \circ X|_U:U \to \mathbb R^{2n}
\]
es suave. 3. Pero por la definición de las funciones $a_i$ tenemos que $X_p=\sum_{i=1}^n a_i(p)\frac{\partial}{\partial x_i}|_p$, osea las coordenadas de $X_p$ son $a_i(p)$. Por lo tanto

\[
\widetilde\varphi \circ X|_U(p)=\widetilde{\varphi}(p,X_{p})=(x_1(p),\ldots,x_n(p),a_1(p),\ldots,a_n(p))
\]
entonces cada funcion coordenada es suave. En particular, las $a_i$ son suaves.

\begin{itemize}
\item $2 \Rightarrow 3$.
\end{itemize}
\begin{enumerate}
\item Sea $V$ abierto de $M$ y sea $f \in C^\infty(V)$. Quiero ver $Xf:V\rightarrow\mathbb{R}$ es suave para esto alcanza con chequear suavidad con cartas suaves de $V$
\item Tomemos una carta $(U,\varphi=(x_1,\ldots,x_n))$ de $V$. Entonces para todo $p\in U$ tenemos
\end{enumerate}
\[
X_p=\sum_{i=1}^n a_i(p)\frac{\partial}{\partial x_i}|_p
\]
, con $a_i$ suaves. 3. Tenemos $Xf:U \to \mathbb R$ con

\[
(Xf)(p)=X_p f=\sum_{i=1}^n a_i(p)\frac{\partial }{\partial x_i}\bigg|_{p}f
\]
\begin{enumerate}[start=4]
\item Como las funciones coordenadas $a_{i}(p)$ son suaves basta ver que
\end{enumerate}
\[
\frac{\partial}{\partial x_{i}} f:U\subseteq M\rightarrow \mathbb{R}\quad\text{ dada por }\quad \frac{\partial}{\partial x_{i}}f(p)=\frac{\partial}{\partial x_{i}}\bigg|_{p}f
\]
son suaves (como funciones que reciben $p\in U$) 5. Pero lo son , por que

\[
\left(\frac{\partial}{\partial x_{i}}f\right)(p)=\frac{\partial}{\partial x_{i}}\bigg|_{p}f=\frac{\partial}{\partial r_{i}}\bigg|_{\varphi(p)}f\circ\varphi ^{-1}=\frac{\partial}{\partial r_{i}} f\circ\varphi ^{-1}(\varphi(p))
\]
osea

\[
\frac{\partial}{\partial x_{i}}f =\left(\frac{\partial}{\partial r_{i}}f\circ\varphi ^{-1}\right)\circ\varphi
\]
pero como $f\circ\varphi ^{-1}$ es suave entonces $\frac{\partial}{\partial r_{i}}f\circ\varphi ^{-1}$ es suave y como $\varphi$ es suave entonces $\frac{\partial}{\partial x_{i}}f$ es suave

\begin{itemize}
\item $3 \Rightarrow 1$.
\end{itemize}
\begin{enumerate}
\item Supongamos que para todo abierto $V\subseteq M$ y para toda función $f\in C^\infty(V)$, se tiene $Xf\in C^\infty(V)$. Queremos ver que $X:M\to TM$ es suave.
\item Para esto tomamos una carta suave $(U,\varphi=(x_1,\ldots,x_n))$ de $M$. Esta induce una carta suave del fibrado tangente
\end{enumerate}
\[
(\widetilde U=\pi^{-1}(U),\widetilde\varphi).
\]
\begin{enumerate}[start=3]
\item Entonces basta ver que
\end{enumerate}
\[
\widetilde\varphi\circ X|_U:U\to\mathbb R^{2n}
\]
es suave. (Esto es por que $p\in U$ y es directo ver que $X|_{U}(U)\subseteq \tilde{U}$ osea que despue es componer a derecha con $\varphi ^{-1}$ que es suave)\\
4. Para cada $p\in U$ tenemos

\[
\widetilde\varphi\circ X|_U(p)=\left(x_1(p),\ldots,x_n(p),a_1(p),\ldots,a_n(p)\right),
\]
donde $a_1(p),\ldots,a_n(p)$ son las coordenadas de $X_p$ en la base coordenada. 5. Las primeras $n$ componentes son suaves, porque $x_1,\ldots,x_n$ son las funciones coordenadas de una carta suave.\\
6. Por lo tanto, falta ver que

\[
a_i:U\to\mathbb R
\]
es suave para cada $i=1,\ldots,n$.\\
7. Para cada $p\in U$, tenemos

\[
X_p=\sum_{i=1}^n a_i(p)\frac{\partial}{\partial x_i}\bigg|_p.
\]
\begin{enumerate}[start=8]
\item Evaluamos en la función coordenada $x_i:U\to\mathbb R$. Entonces
\end{enumerate}
\[
X_px_i=a_i(p).
\]
\begin{enumerate}[start=9]
\item Por lo tanto,
\end{enumerate}
\[
a_i=Xx_i:U\to\mathbb R.
\]
\begin{enumerate}[start=10]
\item Pero $x_i\in C^\infty(U)$, entonces por hipótesis $Xx_i\in C^\infty(U)$. Luego $a_i$ es suave para cada $i$.
\item Así, todas las funciones componentes de $\widetilde\varphi\circ X|_U$ son suaves. Por lo tanto $\widetilde\varphi\circ X|_U$ es suave.
\item Como esto vale para cualquier punto $p$ del fibrado tangente, concluimos que $X:M\to TM$ es suave.
\end{enumerate}
\begin{itemize}
\item $3 \Rightarrow 4$ es inmediato tomando $V=M$.
\item $4 \Rightarrow 3$.
\end{itemize}
\begin{enumerate}
\item Queremos probar que si $V\subseteq M$ es abierto y $f\in C^\infty(V)$, entonces $Xf\in C^\infty(V)$.
\item Lo único que sabemos por hipótesis es que si $g\in C^\infty(M)$, entonces $Xg\in C^\infty(M)$.
\item Sea $p\in V$ arbitrario. Como $V$ es abierto, por \hyperref[blk:405b16]{Extensión de funciones suaves} existe un abierto $U\subseteq M$ tal que
\end{enumerate}
\[
p\in U \qquad \text{y} \qquad \overline{U}\subseteq V.
\]
y una funcion $\widetilde f\in C^\infty(M)$ tal que $\widetilde f|_{ U}=\left.f\right|_{ U}$ 4. Por hipótesis, como $\widetilde f\in C^\infty(M)$, entonces

\[
X\widetilde f\in C^\infty(M).
\]
\begin{enumerate}[start=5]
\item Además, como $\left.\widetilde f\right|_U=\left.f\right|_U$, entonces para todo $q\in U$ las funciones $\widetilde f$ y $f$ coinciden en un abierto alrededor de $q$. Por lo tanto,
\end{enumerate}
\[
X_q\left(\left.\widetilde f\right|_U\right)=X_q\left(\left.f\right|_U\right).
\]
por \hyperref[blk:b713ca]{Localidad de los vectores tangentes} 6. Como esto vale para todo $q$ tenemos que,

\[
(X\widetilde f)|_U=(Xf)|_U.
\]
\begin{enumerate}[start=7]
\item Como $X\widetilde f\in C^\infty(M)$, su restricción a $U$ es suave. Luego
\end{enumerate}
\[
(Xf)|_U\in C^\infty(U).
\]
\begin{enumerate}[start=8]
\item Como $p\in V$ era arbitrario, para cada punto de $V$ existe un abierto $U\subseteq V$ alrededor de ese punto tal que $(Xf)|_U$ es suave. Por lo tanto,
\end{enumerate}
\[
Xf\in C^\infty(V).
\]
notar que usamos \hyperref[blk:e32b2f]{Observacion}

\bookqed
\end{bookproof}
\section{Campos suaves y subvariedades}
\begin{bookproposition}{Suavidad de un campo en una carta}
Sea $(U,\varphi=(x_1,\ldots,x_n))$ una carta suave de $M$ y sea $X:U \to TU$ un campo vectorial. Entonces $X$ es suave si y solo si las coordenadas de $X$ en el marco $\{\frac{\partial}{\partial x_1},\ldots,\frac{\partial}{\partial x_n}\}$ son funciones suaves.

\end{bookproposition}
\begin{bookproof}{}
\begin{itemize}
\item $(\Rightarrow)$ Caso particular de \hyperref[blk:3accf8]{Caracterizaciones de suavidad}
\item $(\Leftarrow)$
\end{itemize}
\begin{enumerate}
\item Tenemos la carta asociada a la carta $\varphi$
\end{enumerate}
\[
\bigg(\widetilde U=\pi^{-1}(U),\tilde\varphi=\left(x_1\circ\pi,\ldots,x_n\circ\pi,v_1,\ldots,v_n\right)\bigg)
\]
que en este caso es un difeo entre $TU$ y $\varphi(U)\times\mathbb{R}^n$.\\
2. Por tanto, $X$ es suave si $\widetilde\varphi\circ X$ es suave (por que si esto es suave componer a derecha con $\varphi ^{-1}:\varphi(U)\rightarrow U$ sigue siendo suave, por que $\varphi ^{-1}$ es homeo, osea que al componer no cambiamos codominio) 3. Como

\[
\widetilde\varphi\circ X:U\subseteq M\longrightarrow\mathbb{R}^{2n}\quad\text{esta dada por}\quad p\longmapsto \left(x_1(p),\ldots,x_n(p),a_1(p),\ldots,a_n(p)\right)
\]
donde $a_i:U\to\mathbb{R}$ son las coordenadas de $X_{p}$ con respecto al marco

\[
\left\{\left.\frac{\partial}{\partial x_1}\right|_p,\ldots,\left.\frac{\partial}{\partial x_n}\right|_p\right\},
\]
\begin{enumerate}[start=4]
\item Y por hipotesis son suaves las coordenadas $a_{i}$ son suaves (obviemente $x_{i}(p)$ son suaves por ser carta), se sigue que $\widetilde\varphi\circ X$ es suave
\item Entonces, para cualquier punto $p$ podemos encontrar cartas que cumplen la definicion, por \hyperref[blk:funcion-suave-entre-variedades]{localidad de la suavidad}.
\end{enumerate}
\bookqed
\end{bookproof}
\section{Curvas integrales}
\begin{booktheorem}{Existencia y unicidad de las curvas integrales}
Sea $X \in \mathfrak X(M)$. Entonces, para cada $p \in M$, existen $a(p)\in\{-\infty\}\cup\mathbb R$, $b(p)\in\mathbb R\cup\{+\infty\}$ y una curva suave $\gamma_p:(a(p),b(p))\subset\mathbb R\to M$ tal que:

\begin{itemize}
\item (a) $0\in(a(p),b(p))$ y $\gamma_p(0)=p$.
\item (b)  $\gamma_p$ es curva integral de $X$: $\gamma_p'(t)=X_{\gamma_p(t)}$.
\item (c) Maximalidad: si $\sigma:(c,d)\subset\mathbb R\to M$ es curva suave que cumple 1 y 2, entonces $(c,d)\subset(a(p),b(p))$ y $\sigma(t)=\gamma_p(t)$ para todo $t\in(c,d)$.
\end{itemize}
En consecuencia, $\gamma_p$ es la unica curva integral que satisface 1, 2 y 3, y es llamada curva integral maximal de $X$ que comienza en $p$.

\end{booktheorem}
\begin{bookproof}{}
\begin{enumerate}
\item Proponemos $(a(p),b(p))$ como la union de todos los intervalos abiertos que contienen al instante $0$ y que son dominios de curvas integrales de $X$ que en el instante $0$ pasan por $p$ (osea que los dominios de estas curvas son intervalos)
\item Para ver que $(a(p),b(p))$ es no vacio, basta exhibir una curva integral que en $t=0$ pasa por $p$ y que este definida en algun intervalo chiquito de $t=0$.
\item Para construir tal curva, basta en hacerlo dentro del dominio de una carta $(U,\varphi=(x_1,\ldots,x_n))$ y bajar el problema a $\mathbb R^n$.
\item Necesitamos $\sigma:(-\epsilon,\epsilon)\to U$ tal que $\sigma(0)=p$ y $\sigma'(t)=X_{\sigma(t)}$.
\item Luego usando la 3era parte de \hyperref[blk:6da36a]{Ejercicio} tenemos que si $\widetilde \sigma=\varphi\circ\sigma$ osea
\end{enumerate}
\[
\widetilde \sigma(t)=(y_1(t),\ldots,y_n(t))
\]
entonces

\[
\sigma'(t)=\sum^{n}_{k=1} y_{k}'(t)\frac{\partial}{\partial x_{k}}\bigg|_{\sigma(k)}
\]
\begin{enumerate}[start=6]
\item Por otro lado como $X$ es suave, podemos usar el marco $\big\{ \frac{\partial}{\partial x_{1}},\ldots \frac{\partial}{\partial x_{n}} \big\}$ para expresar $X|_{U}$ como
\end{enumerate}
\[
X|_U=\sum_{k=1}^n a_k\frac{\partial}{\partial x_k}
\]
con $a_k:U\to\mathbb R$ funciones suaves. 7. Entonces

\[
X_{\sigma(t)}=\sum_{k=1}^n a_k(\sigma(t))\left.\frac{\partial}{\partial x_k}\right|_{\sigma(t)}.
\]
\begin{enumerate}[start=8]
\item Como queremos $\sigma'(t)=X_{\sigma(t)}$ y que $\sigma(0)=p$, usando coordenadas queda el problema
\end{enumerate}
\[
\begin{cases}y_1'(t)=a_1(\sigma(t))=a_1\bigl(\varphi^{-1}(y_1(t),\dots,y_n(t))\bigr)=b_1(y_1(t),\dots,y_n(t)),\\\qquad\vdots\\y_n'(t)=a_n(\sigma(t))=a_n\bigl(\varphi^{-1}(y_1(t),\dots,y_n(t))\bigr)=b_n(y_1(t),\dots,y_n(t)),\\\varphi(p)=(y_1(0),\dots,y_n(0)).\end{cases}
\]
obs: Aca estamos resolviendo el problema en $\varphi(U)$. El campo en $\varphi(U)$ es

\[
\widetilde X_{(r_1,\ldots,r_n)}=\sum_{k=1}^n a_k(\varphi^{-1}(r_1,\ldots,r_n))\left.\frac{\partial}{\partial r_k}\right|_{(r_1,\ldots,r_n)}.
\]
\begin{enumerate}[start=9]
\item por \hyperref[blk:81ef7e]{Fundamental theorem for autonomous ODEs}, tal problema tiene solucion en el instante $t_0=0$ osea existe un intervalo abierto alrededor de $0$, digamos $(-\epsilon,\epsilon)$, y una curva
\end{enumerate}
\[
y:(-\epsilon,\epsilon)\to \widetilde{U}\subset\mathbb R^n
\]
que resuelve el problema 10. Por tanto $\sigma=\varphi^{-1}\circ y$ es curva integral de $X$ que pasa por $p$ en $t=0$. 11. Por lo tanto la union de dominios es no vacio por que $(-\epsilon,\epsilon)$ está en ella 12. Ahora necesitamos definir

\[
\gamma_p:(a(p),b(p))\to M
\]
proponemos lo siguiente: Dado $t\in(a(p),b(p))$ arbitrario. Por definicion de $(a(p),b(p))$, existe una curva integral $\alpha:I\subset\mathbb R\to M$ de $X$ que en $t=0$ pasa por $p$ y tal que $t\in I$ con $I$ intervalo abierto del $0$. Luego podemos definir

\[
\gamma_p(t):=\alpha(t).
\]
\begin{enumerate}[start=13]
\item Debemos ver la buena definicion de $\gamma_p$: Entonces miremos otra curva $\beta:J\subset\mathbb R\to M$ que es otra curva integral de $X$ que en $t=0$ pasa por $p$ y con $t\in J$ con $J$ intervalo abierto del $0$
\item Sea
\end{enumerate}
\[
\Omega=\{s\in I\cap J:\alpha(s)=\beta(s)\}.
\]
Vemos que $\Omega=I\cap J$. 15. Como $I,J$ son intervalos abiertos de $0$, entonces $I\cap J$ es un intervalo abierto de $0$, en particular conexo. 16. Primero notamos que $\Omega$ es no vacio pues $0$ está en él ya que $\alpha (0)=p=\beta(0)$\\
17. \textbf{$\Omega$ es cerrado en $I\cap J$:} En efecto, si $(s_m)$ es una sucesión en $\Omega$ tal que $s_m\to s$ con $s\in I\cap J$, entonces $\alpha(s_m)=\beta(s_m)$ para todo $m$. Como $\alpha$ y $\beta$ son continuas, se tiene $\alpha(s_m)\to\alpha(s)$ y $\beta(s_m)\to\beta(s)$. Por lo tanto, $\alpha(s)=\beta(s)$, es decir, $s\in\Omega$. Luego $\Omega$ es cerrado en $I\cap J$. 18. \textbf{$\Omega$ es abierto en $I\cap J$:} Sea $s_0\in\Omega$. Entonces $\alpha(s_0)=\beta(s_0)=q$. Tomamos una carta $(V,\psi)$ alrededor de $q$. Como $V$ es abierto en $M$ y $\alpha,\beta$ son continuas, el conjunto

\[
K=(I\cap J)\cap\alpha^{-1}(V)\cap\beta^{-1}(V)
\]
es abierto en $I\cap J$ ademas no vacio por que $s_0\in K$, porque $\alpha(s_0)=\beta(s_0)=q\in V$. Por definición de $K$, para todo $s\in K$ se tiene $\alpha(s)\in V$ y $\beta(s)\in V$. Por lo tanto, podemos mirar las curvas coordenadas $\widetilde\alpha=\psi\circ\alpha|_K$ y $\widetilde\beta=\psi\circ\beta|_K$. Como $\alpha$ y $\beta$ son curvas integrales de $X$, las curvas $\widetilde\alpha$ y $\widetilde\beta$ resuelven el mismo sistema local de EDO asociado a $X$ en la carta $(V,\psi)$ (Esto sale de la misma manera que la 1era parte). Además, pasan por el mismo punto en el instante $s_0$, porque

\[
\widetilde\alpha(s_0)=\psi(\alpha(s_0))=\psi(q)=\psi(\beta(s_0))=\widetilde\beta(s_0).
\]
Por unicidad del teorema fundamental de EDO, $\widetilde\alpha$ y $\widetilde\beta$ coinciden en su dominio común, que es $K$. Entonces, para todo $s\in K$, tenemos

\[
\psi(\alpha(s))=\psi(\beta(s)).
\]
Como $\psi$ es inyectiva, se sigue que $\alpha(s)=\beta(s)$ para todo $s\in K$. Por lo tanto, $K\subset\Omega$. Como $K$ es un abierto de $I\cap J$ que contiene a $s_0$, concluimos que $\Omega$ es abierto en $I\cap J$. 19. Asi, $\Omega$ es abierto y cerrado en el conexo $I\cap J$, por tanto $\Omega=I\cap J$. En particular, si $t\in I\cap J$, entonces $\alpha(t)=\beta(t)$. 20. Luego $\gamma_p$ esta bien definida. 21. \textbf{$\gamma_p:(a(p),b(p))\to M$ es suave y es curva integral de $X$:} Sea $t\in(a(p),b(p))$. Por definición de $(a(p),b(p))$, existe una curva integral $\alpha:I\to M$ de $X$, con $0\in I$, $\alpha(0)=p$ y $t\in I$. 22. Como $(a(p),b(p))$ es union de intervalos abiertos, es abierto en $\mathbb R$. Ademas, el intervalo $I$ es uno de los que aparecen en esa union, luego $I\subset(a(p),b(p))$. 23. Afirmamos que $\gamma_p|_I=\alpha$. En efecto, si $s\in I$, entonces por la definición de $\gamma_p(s)$ podemos usar la curva integral $\alpha$ para calcular ese valor. Por lo tanto, $\gamma_p(s)=\alpha(s)$. 24. Como $\alpha$ es suave, se sigue que $\gamma_p|_I$ es suave. 25. Como esto vale para todo $t\in(a(p),b(p))$, la curva $\gamma_p$ es suave. 26. Además, como $\gamma_p|_I=\alpha$ y $\alpha$ es curva integral de $X$, para todo $s\in I$ tenemos $\alpha'(s)=X_{\alpha(s)}$. 27. Por lo tanto, para todo $s\in I$,

\[
\gamma_p'(s)=\alpha'(s)=X_{\alpha(s)}=X_{\gamma_p(s)}.
\]
\begin{enumerate}[start=28]
\item Como $t$ era arbitrario, concluimos que
\end{enumerate}
\[
\gamma_p'(t)=X_{\gamma_p(t)}
\]
para todo $t\in(a(p),b(p))$. Luego $\gamma_p$ es curva integral de $X$. 29. También se tiene $\gamma_p(0)=p$, porque si tomamos cualquier curva integral $\alpha:I\to M$ usada en la construcción, entonces $\gamma_p(0)=\alpha(0)=p$. 30. \textbf{Maximalidad:} Sea $\sigma:(c,d)\to M$ una curva integral de $X$ tal que $0\in(c,d)$ y $\sigma(0)=p$. 31. Entonces $(c,d)$ es uno de los intervalos que aparecen en la unión que define $(a(p),b(p))$. Por lo tanto,

\[
(c,d)\subset(a(p),b(p)).
\]
\begin{enumerate}[start=32]
\item Además, para todo $t\in(c,d)$, la curva $\sigma$ puede usarse en la definición de $\gamma_p(t)$. Por lo tanto,
\end{enumerate}
\[
\gamma_p(t)=\sigma(t).
\]
\begin{enumerate}[start=33]
\item Así, toda curva integral de $X$ que pasa por $p$ en el instante $0$ está contenida en $\gamma_p$ y coincide con $\gamma_p$ en su dominio.
\item Por lo tanto, $\gamma_p$ es maximal.
\item Finalmente, si hubiese otra curva maximal $\eta:(A,B)\to M$ con $\eta(0)=p$ y $\eta'(t)=X_{\eta(t)}$, entonces por la maximalidad de $\gamma_p$ tendríamos $(A,B)\subset(a(p),b(p))$ y $\eta=\gamma_p$ en $(A,B)$.
\item Por la maximalidad de $\eta$, también tendríamos $(a(p),b(p))\subset(A,B)$ y $\gamma_p=\eta$ en $(a(p),b(p))$.
\item Luego
\end{enumerate}
\[
(A,B)=(a(p),b(p))
\]
y

\[
\eta=\gamma_p.
\]
Por lo tanto, la curva integral maximal que comienza en $p$ es única.

\bookqed
\end{bookproof}
\begin{booklemma}{}
Sean $M$ una variedad suave, $X\in\mathfrak{X}(M)$ y

\[
\gamma_p:(a(p),b(p))\longrightarrow M
\]
la curva integral maximal de $X$ que comienza en $p$, es decir, $\gamma_p(0)=p$. Si

\[
q=\gamma_p(t_0),
\]
entonces la curva integral maximal que comienza en $q$ está dada por

\[
\gamma_q(t)=\gamma_p(t+t_0),
\]
y su intervalo maximal es

\[
(a(q),b(q))=(a(p)-t_0,b(p)-t_0).
\]
\end{booklemma}
\begin{bookproof}{}
38. Sea

\[
I_p=(a(p),b(p))
\]
y consideremos el intervalo trasladado

\[
J=I_p-t_0=(a(p)-t_0,b(p)-t_0).
\]
\begin{enumerate}[start=39]
\item Definimos
\end{enumerate}
\[
\beta:J\longrightarrow M,\qquad \beta(t)=\gamma_p(t+t_0).
\]
\begin{enumerate}[start=40]
\item En primer lugar,
\end{enumerate}
\[
\beta(0)=\gamma_p(t_0)=q.
\]
\begin{enumerate}[start=41]
\item Además, como $\gamma_p$ es una curva integral de $X$,
\end{enumerate}
\[
\beta'(t)=\frac{d}{dt}\gamma_p(t+t_0)=X_{\gamma_p(t+t_0)}=X_{\beta(t)}.
\]
Por lo tanto, $\beta$ es una curva integral de $X$ que comienza en $q$.

\begin{enumerate}[start=42]
\item Veamos que $\beta$ es maximal. Supongamos que admite una extensión integral
\end{enumerate}
\[
\widetilde{\beta}:\widetilde{J}\longrightarrow M,\qquad J\subsetneq\widetilde{J}.
\]
\begin{enumerate}[start=43]
\item Definimos entonces
\end{enumerate}
\[
\widetilde{\gamma}(s)=\widetilde{\beta}(s-t_0),\qquad s\in\widetilde{J}+t_0.
\]
Esta es una curva integral de $X$.

\begin{enumerate}[start=44]
\item Para todo $s\in I_p$,
\end{enumerate}
\[
\widetilde{\gamma}(s)=\widetilde{\beta}(s-t_0)=\beta(s-t_0)=\gamma_p(s).
\]
\begin{enumerate}[start=45]
\item Así, $\widetilde{\gamma}$ extiende a $\gamma_p$, pues
\end{enumerate}
\[
I_p=J+t_0\subsetneq\widetilde{J}+t_0,
\]
lo cual contradice la maximalidad de $\gamma_p$.

\begin{enumerate}[start=46]
\item Por consiguiente, $\beta$ es la curva integral maximal que comienza en $q$. Por unicidad,
\end{enumerate}
\[
\gamma_q=\beta.
\]
\begin{enumerate}[start=47]
\item En conclusión,
\end{enumerate}
\[
\gamma_q(t)=\gamma_p(t+t_0)
\]
y

\[
(a(q),b(q))=(a(p)-t_0,b(p)-t_0).
\]
\bookqed
\end{bookproof}
\section{Groundwork de generador infinitesimal}
\begin{booktheorem}{}
Sea $X\in\mathfrak X(M)$ y sea $\theta:\mathcal D\subseteq\mathbb R\times M\to M$ su flujo. Entonces:

\begin{itemize}
\item \textbf{(i)} $\bigcup_{t>0}D_t=M$.
\item \textbf{(ii)} Dados $s,t\in\mathbb R$, entonces $\operatorname{Dom}(\theta_t\circ\theta_s)\subseteq D_{t+s}$. Si $s$ y $t$ tienen el mismo signo, entonces $\operatorname{Dom}(\theta_t\circ\theta_s)=D_{t+s}$.
\item \textbf{(iii)} Para todo $p\in\operatorname{Dom}(\theta_s\circ\theta_t)$ se cumple
\end{itemize}
\[
\theta_s\circ\theta_t(p)=\theta_{s+t}(p).
\]
\begin{itemize}
\item \textbf{(iv)} La función
\end{itemize}
\[
\theta_t:D_t\to D_{-t}
\]
es un difeomorfismo, con inversa $\theta_{-t}$.

\end{booktheorem}
\begin{bookproof}{}
\begin{itemize}
\item \textbf{(i)}
\end{itemize}
\begin{enumerate}
\item Dado $p\in M$ arbitrario, por \hyperref[blk:3ff110]{Existencia y unicidad de las curvas integrales} sabemos que existe una curva integral que comienza en $p$ definida en algún intervalo maximal $(a(p),b(p))$.
\item Como dicho intervalo contiene al $0$. Tomamos algun $t\in(0,\epsilon)\subseteq (a(p),b(p))$. Entonces $p\in D_t$
\item Como $p$ era arbitrario, se concluye que todo punto de $M$ pertenece a algún $D_t$ con $t>0$.
\end{enumerate}
\[
\bigcup_{t>0}D_t=M
\]
(Notar que claramente $D_{t}\subseteq M$ para cualquier $t$ por eso vale trivialmente la otra inclusion)

\begin{itemize}
\item \textbf{(ii)}
\end{itemize}
\begin{enumerate}
\item Sea $p\in\operatorname{Dom}(\theta_t\circ\theta_s)$. Esto significa que $p\in D_s$
\item Luego esta bien definida $\theta_{s}(p)$ y ademas $q=\theta_s(p)\in D_t$.
\item Luego por definicion $q=\gamma_p(s)$
\item Por la observación sobre traslación de curvas integrales maximales, la curva integral maximal que comienza en $q$ satisface
\end{enumerate}
\[
(a(q),b(q))=(a(p)-s,b(p)-s)
\]
\begin{enumerate}[start=5]
\item Como $q\in D_t$, se tiene $t\in(a(q),b(q))$.
\item Usando la igualdad de intervalos, esto equivale a
\end{enumerate}
\[
t+s\in(a(p),b(p))
\]
\begin{enumerate}[start=7]
\item Por lo tanto $p\in D_{t+s}$, y se obtiene
\end{enumerate}
\[
\operatorname{Dom}(\theta_t\circ\theta_s)\subseteq D_{t+s}
\]
\begin{enumerate}[start=8]
\item Ahora supongamos que $s,t>0$ y sea $p\in D_{s+t}$.
\item Entonces $s+t\in(a(p),b(p))$ y, como $0<s<s+t$, también $s\in(a(p),b(p))$.
\end{enumerate}
\[
p\in D_s
\]
\begin{enumerate}[start=10]
\item Si $q=\theta_s(p)=\gamma_p(s)$, entonces la curva integral maximal que comienza en $q$ tiene dominio
\end{enumerate}
\[
(a(q),b(q))=(a(p)-s,b(p)-s)
\]
\begin{enumerate}[start=11]
\item Como $s+t\in(a(p),b(p))$, se obtiene $t\in(a(q),b(q))$.
\end{enumerate}
\[
q\in D_t
\]
\begin{enumerate}[start=12]
\item Luego $p\in\operatorname{Dom}(\Theta_t\circ\Theta_s)$.
\item Por lo tanto, si $s,t>0$, se cumple
\end{enumerate}
\[
\operatorname{Dom}(\theta_t\circ\theta_s)=D_{t+s}
\]
\begin{enumerate}[start=14]
\item El caso $s,t<0$ se prueba de manera análoga, usando que $s+t<s<0$ y la misma traslación del dominio maximal.
\end{enumerate}
\begin{itemize}
\item \textbf{(iii)}
\end{itemize}
\begin{enumerate}
\item Sea $p\in\operatorname{Dom}(\theta_s\circ\theta_t)$ y sea $q=\theta_t(p)$.
\item Entonces $q=\gamma_p(t)$, y por la descripción de la curva maximal que comienza en $q$,
\end{enumerate}
\[
\gamma_q(s)=\gamma_p(t+s)
\]
\begin{enumerate}[start=3]
\item Por lo tanto,
\end{enumerate}
\[
\theta_s(\theta_t(p))=\theta_s(q)=\gamma_q(s)=\gamma_p(t+s)=\theta_{s+t}(p)
\]
\begin{enumerate}[start=4]
\item Concluimos que
\end{enumerate}
\[
\theta_s\circ\theta_t(p)=\theta_{s+t}(p)
\]
\begin{itemize}
\item \textbf{(iv)}
\end{itemize}
\begin{enumerate}
\item Veamos primero que $\theta_t(D_t)\subseteq D_{-t}$.
\item Sea $p\in D_t$ y sea $q=\theta_t(p)=\gamma_p(t)$.
\item Nuevamente,
\end{enumerate}
\[
(a(q),b(q))=(a(p)-t,b(p)-t)
\]
\begin{enumerate}[start=4]
\item Como $0\in(a(p),b(p))$, se tiene $-t\in(a(q),b(q))$.
\end{enumerate}
\[
q\in D_{-t}
\]
\begin{enumerate}[start=5]
\item Entonces $\theta_t:D_t\to D_{-t}$ está bien definida.
\item Además, usando la parte \textbf{(iii)},
\end{enumerate}
\[
\theta_{-t}(\theta_t(p))=\theta_0(p)=\gamma_{p}(0)=p
\]
\begin{enumerate}[start=7]
\item Análogamente, para $q\in D_{-t}$,
\end{enumerate}
\[
\theta_t(\theta_{-t}(q))=\theta_0(q)=\gamma_{q}(0)=q
\]
\begin{enumerate}[start=8]
\item Así, $\theta_{-t}$ es la inversa de $\theta_t$.
\item Veamos que \textbf{$\Theta_{t}$ es suave}
\item Fijemos $t\in\mathbb R$. Recordemos que el flujo es una aplicación suave
\end{enumerate}
\[
\Theta:\mathcal D\subseteq\mathbb R\times M\to M
\]
\begin{enumerate}[start=11]
\item Por definición,
\end{enumerate}
\[
D_t=\{p\in M:(t,p)\in\mathcal D\}
\]
y $\Theta_t:D_t\to M$ está dada por $\Theta_t(p)=\Theta(t,p)$ 12. Primero observemos que $D_t$ es abierto en $M$. En efecto, si definimos

\[
\iota_t:M\to\mathbb R\times M
\]
por $\iota_t(p)=(t,p)$, entonces $\iota_t$ es continua y $D_t=\iota_t^{-1}(\mathcal D)$. 13. Como $\mathcal D$ es abierto, se sigue que $D_t$ es abierto. Con lo cual es subvariedad de $M$ 14. Ahora consideramos la restricción

\[
\iota_t|_{D_t}:D_t\to\mathcal D
\]
dada por $\iota_t|_{D_t}(p)=(t,p)$. Esta aplicación es suave porque es la restricción de una aplicación suave a un abierto dentro de la variedad. 15. Para todo $p\in D_t$, tenemos

\[
\Theta_t(p)=\Theta(t,p)=\Theta\big(\iota_t|_{D_t}(p)\big).
\]
\begin{enumerate}[start=16]
\item Por lo tanto,
\end{enumerate}
\[
\Theta_t=\Theta\circ \iota_t|_{D_t}.
\]
\begin{enumerate}[start=17]
\item Como $\iota_t|_{D_t}$ es suave y $\Theta$ es suave, se concluye que $\Theta_t:D_t\to M$ es suave.
\item Además, ya se probó que $\Theta_t(D_t)\subseteq D_{-t}$. Como $D_{-t}$ es un abierto de $M$, la misma aplicación puede verse como
\end{enumerate}
\[
\Theta_t:D_t\to D_{-t}.
\]
\begin{enumerate}[start=19]
\item Esta aplicación también es suave como aplicación con codominio $D_{-t}$, porque la suavidad hacia un abierto de $M$ se verifica componiendo con la inclusión $D_{-t}\hookrightarrow M$. Y usando continuidad de $\Theta_{t}$. Además, $D_{-t}$ es una subvariedad incrustada de $M$ por ser un abierto de $M$.
\item Por lo tanto, $\Theta_t:D_t\to D_{-t}$ es un difeomorfismo con inversa $\Theta_{-t}$.
\end{enumerate}
\bookqed
\end{bookproof}
\begin{booklemma}{Uniformidad global del tiempo}
Sea $X$ un campo suave sobre una variedad $M$ y sea $\Theta$ su flujo. Supongamos que existe un número positivo $\varepsilon>0$ tal que, para todo $p\in M$, el dominio de $\gamma_p$ contiene a $(-\varepsilon,\varepsilon)$.

Entonces $X$ es un campo completo.

\end{booklemma}
\begin{bookproof}{}
\begin{enumerate}
\item La idea es que, cuando se intenta acabar una curva integral maximal que comienza en $p$, la estiramos antes de que termine usando el intervalo uniforme $(-\varepsilon,\varepsilon)$, que sirve para todo punto.
\item Razonemos por el absurdo. Supongamos que existe $p\in M$ tal que $b(p)<+\infty$. Osea estamos suponiendo que $X$ no es campo completo. (El caso $a(p)>-\infty$ es análogo)
\item Elegimos un punto del final de la curva, antes de que se acabe, faltando menos de $\varepsilon$ segundos.
\item Tomamos $t_0\in(a(p),b(p))$ tal que
\end{enumerate}
\[
b(p)-\varepsilon<t_0<b(p)
\]
\begin{enumerate}[start=5]
\item Pensemos en la curva integral maximal de $q$
\end{enumerate}
\[
q=\gamma_p(t_0)
\]
por hipótesis, está definida al menos en $(-\varepsilon,\varepsilon)$. 6. Definimos una curva $\alpha$ pegando $\gamma_p$ con la curva que empieza en $q$:

\[
\alpha(t)=\begin{cases}\gamma_p(t),&a(p)<t<b(p),\\\gamma_q(t-t_0),&-\varepsilon<t-t_0<\varepsilon.\end{cases}
\]
\begin{enumerate}[start=7]
\item Primero veamos que está bien definida donde se solapan las dos definiciones.
\item Notar que $(a(p),b(p))\cap(t_0-\varepsilon,t_0+\varepsilon)=(t_0-\varepsilon,b(p))$.
\item En efecto, como $a(p)<b(p)-\varepsilon$, podemos elegir $t_0$ de modo que $a(p)<t_0-\varepsilon$. Entonces $a(p)<t_0-\varepsilon<b(p)$.
\item Por otro lado, como $b(p)-\varepsilon<t_0$, se sigue que $b(p)<t_0+\varepsilon$.
\item Por lo tanto, $(a(p),b(p))\cap(t_0-\varepsilon,t_0+\varepsilon)=(t_0-\varepsilon,b(p))$.
\item Entonces n el solapamiento osea si $t\in(t_0-\varepsilon,b(p))$, tenemos que
\end{enumerate}
\[
\gamma_q(t-t_0)=\Theta_{t-t_0}(q)=\Theta_{t-t_0}(\Theta_{t_0}(p))=\Theta_t(p)=\gamma_p(t)
\]
por lo tanto, las dos expresiones coinciden en el solapamiento. 13. Además, $\alpha$ es curva integral de $X$ y en $t=0$ pasa por $p$. 14. Como $b(p)-\varepsilon<t_0$, tenemos $b(p)<t_0+\varepsilon$. 15. Entonces $\alpha$ está definida más allá de $b(p)$. 16. Esto contradice la maximalidad de $\gamma_p$. 17. Luego no puede ocurrir que $b(p)<+\infty$. (Análogamente, no puede ocurrir que $a(p)>-\infty$) 18. Por lo tanto, para todo $p\in M$ se cumple

\[
(a(p),b(p))=\mathbb R
\]
\begin{enumerate}[start=19]
\item Concluimos que $X$ es completo.
\end{enumerate}
\bookqed
\end{bookproof}
\begin{booktheorem}{Campo con soporte compacto}
Sea $X\in\mathfrak X(M)$ con soporte compacto. Entonces $X$ es completo.

Recordemos que

\[
\operatorname{Sop}X=\operatorname{Cl}_M\{q\in M:X_q\neq0\}
\]
\end{booktheorem}
\begin{bookproof}{}
20. Sea $K=\operatorname{Sop}X$. 21. Notemos primero que si $p\notin K$, entonces existe un abierto $V_p$ de $M$ con $p\in V_p$ donde el campo vale cero. 22. Es decir, para todo $q\in V_p$,

\[
X_q=0
\]
\begin{enumerate}[start=23]
\item Entonces, para todo $q\in V_p$, la curva constante
\end{enumerate}
\[
\gamma_q:\mathbb R\to M,\qquad \gamma_q(t)=q
\]
es curva integral maximal de $X$. 24. Por lo tanto, los puntos fuera de $K$ no generan ningún problema: sus curvas integrales están definidas para todo tiempo mientras permanezcan fuera del soporte, y en particular tienen tiempo local uniforme. 25. Falta ver qué pasa en $K$. 26. Por \hyperref[blk:3b18e6]{Uniformidad local del tiempo}, para todo $p\in K$ existe $\varepsilon_p>0$ y existe un abierto $V_p$ de $M$ con $p\in V_p$ tal que, para todo $q\in V_p$, la curva integral que comienza en $q$ está definida en $(-\varepsilon_p,\varepsilon_p)$. Osea $(-\epsilon_{p},\epsilon_{p})\subseteq (a(q),b(q))$ para todo $q\in V_{p}$\\
27. Con los abiertos $V_p$ podemos cubrir a $K$. 28. Como $K$ es compacto, existen $p_1,\ldots,p_s\in K$ tales que

\[
K\subseteq V_{p_1}\cup\cdots\cup V_{p_s}
\]
\begin{enumerate}[start=29]
\item Tomamos
\end{enumerate}
\[
\varepsilon=\min\{\varepsilon_{p_1},\ldots,\varepsilon_{p_s}\}
\]
\begin{enumerate}[start=30]
\item Entonces, para todo $q\in K$, la curva integral que comienza en $q$ está definida al menos en $(-\varepsilon,\varepsilon)$.
\item Por lo tanto, para todo $q\in M$, el dominio de $\gamma_q$ contiene a $(-\varepsilon,\varepsilon)$.
\item Por el lema de uniformidad global del tiempo, $X$ es completo.
\end{enumerate}
\bookqed
\end{bookproof}
\section{Extensión local de campos}
\begin{booktheorem}{Extensión local de campos}
Sea $f:M\to N$ una inmersión y sea $X\in\mathfrak X(M)$. Entonces, para todo $p\in M$, existen un abierto $U$ de $p$ en $M$, un abierto $V$ de $N$ en $f(p)$, con $f(U)\subseteq V$, y un campo suave $\widetilde X\in\mathfrak X(V)$ tal que

\[
\widetilde X_{f(q)}=(df)_qX_q,\qquad q\in U.
\]
Se dice que $\widetilde X$ es una extensión local de $df(X)$.

\end{booktheorem}
\begin{bookproof}{}
\begin{enumerate}
\item Recordemos que, por la forma local de una inmersión, existen cartas cúbicas $(U,\varphi=(x_1,\ldots,x_m))$ centrada en $p$ y $(V,\psi=(y_1,\ldots,y_n))$ centrada en $f(p)$ tales que
\end{enumerate}
\[
\psi\circ f\circ\varphi^{-1}(x_1,\ldots,x_m)=(x_1,\ldots,x_m,0,\ldots,0).
\]
\begin{enumerate}[start=2]
\item Escribimos el campo $X$ en la carta de $M$ como
\end{enumerate}
\[
X|_U=\sum_{k=1}^m a_k\frac{\partial}{\partial x_k},\qquad a_k\in C^\infty(U).
\]
\begin{enumerate}[start=3]
\item Si $q\in U$, entonces (ya lo probamos usando forma local \hyperref[blk:b71da3]{\textasciicircum{}b71da3})
\end{enumerate}
\[
(df)_q\left(\left.\frac{\partial}{\partial x_k}\right|_q\right)=\left.\frac{\partial}{\partial y_k}\right|_{f(q)},\qquad 1\le k\le m.
\]
\begin{enumerate}[start=4]
\item Por lo tanto,
\end{enumerate}
\[
(df)_qX_q=\sum_{k=1}^m a_k(q)\left.\frac{\partial}{\partial y_k}\right|_{f(q)}
\]
\begin{enumerate}[start=5]
\item 
\end{enumerate}
\[
A_k=\begin{cases}a_k\circ\varphi^{-1}\circ\pi\circ\psi,&1\le k\le m,\\0,&m<k\le n,\end{cases}
\]
Esta esta bien definido por que $\psi$ y $\varphi$ son cubicas. 6. Y notamos $A_{k}$ cumple que

\[
(A_{k}\circ f)(p)=a_{k}\circ\varphi ^{-1}\circ\pi\circ\psi\circ f(\varphi ^{-1} \circ\varphi(p))=a_{k}\circ\varphi ^{-1}\circ\varphi(p)=a_{k}(p)
\]
\begin{enumerate}[start=7]
\item Finalmente proponemos
\end{enumerate}
\[
\widetilde X=\sum_{k=1}^n A_k\frac{\partial}{\partial y_k}.
\]
el cual cumple $\widetilde{X}(f(q))=(df)_{q}X_{q}$ para todo $q\in U$ y $A_k$ esta definida en todo $V$, entonces $\widetilde X\in\mathfrak X(V)$ como queriamos.

\bookqed
\end{bookproof}
\begin{booktheorem}{Teorema de enderezamiento}
Sea $X\in\mathfrak X(M)$ y sea $p\in M$ tal que $X_p\neq0$. Entonces existe una carta suave $(U,\varphi=(x_1,\ldots,x_n))$ centrada en $p$ tal que

\[
X|_U=\frac{\partial}{\partial x_1}.
\]
Es decir, en una carta adecuada, el campo se ve como un campo constante.

\end{booktheorem}
\begin{bookproof}{}
8. Por el lema anterior, tomamos una carta $(V,\psi=(y_1,\ldots,y_n))$ centrada en $p$ tal que

\[
\left.\frac{\partial}{\partial y_1}\right|_p=X_p.
\]
\begin{enumerate}[start=9]
\item Por otro lado por \hyperref[blk:3b18e6]{Uniformidad local del tiempo} existe $\widetilde{V}$ abierto de $M$ en $p$ y $\epsilon>0$ tal que el flujo $\Theta$ de $X$ esta definido en $(-\epsilon,\epsilon)\times \widetilde{V}$
\item Achicando si hace falta podes suponer que $\widetilde{V}=V$ y que $(V,\psi)$ es una carta cubica centrada en $p$ osea
\end{enumerate}
\[
\psi(V)=C_{\epsilon}^{n} (0)=(-\epsilon,\epsilon)^{n}
\]
\begin{enumerate}[start=11]
\item Ahora nos vamos a enfocar en una cierta rebanada de $\psi(V)$ dada por
\end{enumerate}
\[
\{ (r_{1},\ldots,r_{n})\in \psi(V):r_{1}=0 \}
\]
que en este caso es igual a $(-\epsilon,\epsilon)^{n-1}$ por ser $\psi$ carta cubica 12. Proponemos

\[
\sigma:(-\epsilon,\epsilon)\times C_{\epsilon}^{n-1}(0)\rightarrow M\qquad (r_{1},(0,r_{2},\ldots r_{n}))\mapsto \Theta_{r_{1}}(\psi ^{-1}(0,r_{2},\ldots,r_{n}))=\Theta_{r_{1}}\circ\psi ^{-1}(0,r_{2},\ldots,r_{n})
\]
\begin{enumerate}[start=13]
\item Ahora vamos a notar que $\sigma$ nos va a ayudar a armar una carta, para esto veamos que $(d\sigma)_{0}$ es invertible. Por que si sucediera entonces es isomorfismo luego por \hyperref[blk:b664ba]{Teorema de la funcion inversa en variedades} tendremos un abierto $U$ abierto del $0$ tal que
\end{enumerate}
\[
\sigma|_{U}^{-1}:W\subseteq M\rightarrow U
\]
sera difeo (osea carta) 14. Luego tenemos si definimos la curva $\alpha (t)=(t,0,\ldots,0)$

\[
\begin{aligned}(d\sigma)_0\left(\frac{\partial}{\partial r_1}\bigg|_0\right)&=(\sigma\circ\alpha )'(0)\\&=\frac{d}{dt}\bigg|_0\sigma(t,0,\dots,0)\\&=\frac{d}{dt}\bigg|_0\theta_t(\psi^{-1}(0,\dots,0))\\&=\frac{d}{dt}\bigg|_0\theta_t(p)\\&=\frac{d}{dt}\bigg|_0\gamma_{p}(t)\\&=\gamma'_{p}(0)\\&=X_{p}\\&=\left.\frac{\partial}{\partial y_1}\right|_p\end{aligned}
\]
en el primer paso usamos \textasciicircum{}99bdaa 15. Por otro lado para $k>1$

\[
\begin{aligned}(d\sigma)_0\frac{\partial}{\partial r_k}\bigg|_0 &=\frac{d}{dt}\bigg|_0\sigma(0,\ldots,t,\ldots,0)\\ &=\frac{d}{dt}\bigg|_0\Theta_{0}\left(\psi^{-1}(0,\ldots,t,\ldots,0)\right)\\ &=\frac{d}{dt}\bigg|_0\psi^{-1}(0,\ldots,t,\ldots,0)\\ &=(d\psi^{-1})_0\frac{\partial}{\partial r_k}\bigg|_0\\ &=\left.\frac{\partial}{\partial y_k}\right|_p. \end{aligned}
\]
el primer igual sale como en \textasciicircum{}99bdaa cambiando la curva $\alpha(t)=te_{k}$. El ultimo igual vale porque $\psi$ es difeomorfismo de carta, entonces $(d\psi^{-1})_0$ es isomorfismo. 16. Entonces $d\sigma$ es isomorfismo. Por tanto, por el teorema de la función inversa existe $U$ abierto de $M$ en $p$ y $W$ abierto de $\mathbb R^n$ en $(0,\dots,0)$ tal que $\sigma:W\to U$ es un difeomorfismo. 17. Ahora nuestra candidata es la carta suave $(U,\phi=\sigma^{-1}|_U=(x_1,\dots,x_n)).$ 18. Veamos que

\[
X|_U=\frac{\partial}{\partial x_1}.
\]
\begin{enumerate}[start=19]
\item Tomamos $q=\sigma(r_1,\ldots,r_n)=\Theta_{r_1}\left(\psi^{-1}(0,r_2,\ldots,r_n)\right).$ Como $\phi=\sigma^{-1}$, osea $\phi(q)=(r_{1},\ldots,r_{n})$ por lo tanto $x_{i}(q)=r_{i}$. Entonces
\end{enumerate}
\[
\begin{aligned}\frac{\partial}{\partial x_1}\bigg|_q&=(d\phi ^{-1})_{\phi(q)}\left(\frac{\partial}{\partial r_{1}}\bigg|_{\phi(q)}\right)\\&=(d\sigma)_{(r_1,\ldots,r_n)}\left(\frac{\partial}{\partial r_1}\bigg|_{(r_1,\ldots,r_n)}\right)\\&=\frac{d}{ds}\bigg|_{s=0}\sigma(r_1+s,r_2,\ldots,r_n)\\&=\frac{d}{ds}\bigg|_{s=0}\Theta_{r_1+s}\left(\psi^{-1}(0,r_2,\ldots,r_n)\right)\\&=\frac{d}{ds}\bigg|_{s=0}\Theta_s\left(\Theta_{r_1}\left(\psi^{-1}(0,r_2,\ldots,r_n)\right)\right)\\&=\frac{d}{ds}\bigg|_{s=0}\Theta_s(q)\\&=\frac{d}{ds}\bigg|_{s=0}\gamma_q(s)\\&=\gamma_q'(0)\\&=X_{\gamma_q(0)}\\&=X_q.\end{aligned}
\]
\begin{enumerate}[start=20]
\item Por lo tanto, para todo $q\in U$, se cumple
\end{enumerate}
\[
\frac{\partial}{\partial x_1}\bigg|_q=X_q.
\]
\begin{enumerate}[start=21]
\item Luego
\end{enumerate}
\[
X|_U=\frac{\partial}{\partial x_1}.
\]
\bookqed
\end{bookproof}
\section{Corchete de Lie}
\begin{bookproposition}{Buena definicion corchete de Lie}
El corchete de Lie está bien definido, es decir, $[X,Y]_p\in T_pM$ y ademas $[X,Y]$ es un campo suave.

\end{bookproposition}
\begin{bookproof}{}
\begin{enumerate}
\item Necesitamos ver que $[X,Y]_p$ es vector tangente.
\item \textbf{Linealidad} Sea $f,g\in C^\infty(M)$. Entonces
\end{enumerate}
\[
[X,Y]_p(f+g)=X_p(Y(f+g))-Y_p(X(f+g))=[X,Y]_pf+[X,Y]_pg,
\]
por linealidad de $X_p$ e $Y_p$. Para ver la igualdad intermedia, por ejemplo,

\[
(Y(f+g))(q)=Y_q(f+g)=Y_qf+Y_qg=(Yf)(q)+(Yg)(q),
\]
y por lo tanto $Y(f+g)=Yf+Yg$ como funciones. Análogo con $X(f+g)$. 3. \textbf{Leibniz} Primero por Leibniz

\[
Y(fg)=fYg+gYf.
\]
Esto se prueba facilmente de la misma manera que en 2. 4. Entonces

\[
X(Y(fg))=X(fYg+gYf)=X(fYg)+X(gYf).
\]
\begin{enumerate}[start=5]
\item Usando Leibniz de $X$,
\end{enumerate}
\[
X(fYg)=X(f)Yg+fX(Yg),
\]
y

\[
X(gYf)=X(g)Yf+gX(Yf).
\]
\begin{enumerate}[start=6]
\item Luego
\end{enumerate}
\[
X(Y(fg))=X(f)Yg+fX(Yg)+X(g)Yf+gX(Yf).
\]
\begin{enumerate}[start=7]
\item Del mismo modo,
\end{enumerate}
\[
X(fg)=fXg+gXf,
\]
y por lo tanto

\[
Y(X(fg))=Y(fXg+gXf)=Y(f)Xg+fY(Xg)+Y(g)Xf+gY(Xf).
\]
\begin{enumerate}[start=8]
\item Restando,
\end{enumerate}
\[
[X,Y](fg)=X(Y(fg))-Y(X(fg)).
\]
\begin{enumerate}[start=9]
\item Los términos cruzados se cancelan porque son productos de funciones reales:
\end{enumerate}
\[
X(f)Yg-Y(g)Xf\equiv0,\qquad X(g)Yf-Y(f)Xg\equiv0
\]
cuando se reordenan los productos entre funciones correspondientes. 10. Queda

\[
[X,Y](fg)=f[X,Y](g)+g[X,Y](f).
\]
entonces

\[
[X,Y]_{p}(fg)=f(g)[X,Y]_{p}(g)+g(p)[X,Y]_{p}(f)
\]
\begin{enumerate}[start=11]
\item Finalmente vimos $[X,Y]_p$ es una vector tangente en $p$, es decir $[X,Y]_p\in T_pM$.
\item \textbf{Suavidad} Para esto usemos \hyperref[blk:3accf8]{Caracterizaciones de suavidad} 4.
\item Entonces dada $f\in C^{\infty}(M)$ tenemos que ver que $[X,Y]f=X(Yf)-Y(Xf)$ es suave
\item Primero notamos $Yf,Xf\in C^{\infty}(M)$ por que $X,Y$ son campos suaves entonces $X(Yf),Y(Xf)\in C^{\infty}(M)$ misma razon
\item Luego $[X,Y]f\in C^{\infty}(M)$ por ser resta de cosas en $C^{\infty}(M)$ por lo tanto $[X,Y]$ es campo suave
\end{enumerate}
\bookqed
\end{bookproof}
\begin{bookproposition}{Propiedades del corchete de Lie}
Sean $X,Y,Z\in\mathfrak X(M)$ y $f,g\in C^\infty(M)$. Entonces:

\begin{itemize}
\item \textbf{(i)} $[X,Y]=-[Y,X]$.
\item \textbf{(ii)} Si $(U,\varphi=(x_{1},\ldots,x_{n}))$ entonces
\end{itemize}
\[
\left[\frac{\partial}{\partial x_i},\frac{\partial}{\partial x_j}\right]=0.
\]
\begin{itemize}
\item \textbf{(iii)} $[X,gY]=g[X,Y]+X(g)Y$.
\item \textbf{(iv)} $[fX,gY]=fg[X,Y]+fX(g)Y-gY(f)X$.
\item \textbf{(v)} Identidad de Jacobi:
\end{itemize}
\[
[X,[Y,Z]]+[Y,[Z,X]]+[Z,[X,Y]]=0.
\]
\end{bookproposition}
\begin{bookproof}{}
\begin{itemize}
\item \textbf{(i)}
\end{itemize}
\begin{enumerate}
\item Sale directo:
\end{enumerate}
\[
-[Y,X]f=-Y(Xf)+X(Yf)=[X,Y]f.
\]
\begin{itemize}
\item \textbf{(ii)}
\end{itemize}
\begin{enumerate}
\item Tomemos $f\in C^\infty(U)$.
\end{enumerate}
\[
\begin{aligned}\left.\frac{\partial}{\partial x_i}\right|_p\left(\frac{\partial f}{\partial x_j}\right)&=\frac{\partial}{\partial r_i}\left[\left(\frac{\partial f}{\partial x_j}\right)\circ\varphi^{-1}\right](\varphi(p))\\&=\frac{\partial}{\partial r_i}\left[\left(\left(\frac{\partial}{\partial r_j}(f\circ\varphi^{-1})\right)\circ\varphi\right)\circ\varphi^{-1}\right](\varphi(p))\\&=\frac{\partial}{\partial r_i}\left[\frac{\partial}{\partial r_j}(f\circ\varphi^{-1})\circ(\varphi\circ\varphi^{-1})\right](\varphi(p))\\&=\frac{\partial}{\partial r_i}\left[\frac{\partial}{\partial r_j}(f\circ\varphi^{-1})\right](\varphi(p))\\&=\frac{\partial^2(f\circ\varphi^{-1})}{\partial r_i\partial r_j}(\varphi(p)).\end{aligned}
\]
\begin{enumerate}[start=2]
\item Y se termina con el teorema de igualdad de las derivadas cruzadas.
\end{enumerate}
\begin{itemize}
\item \textbf{(iii)}
\end{itemize}
\begin{enumerate}
\item Calculamos:
\end{enumerate}
\[
[X,gY]h=X(gYh)-gY(Xh).
\]
\begin{enumerate}[start=2]
\item Usando Leibniz:
\end{enumerate}
\[
X(gYh)=X(g)Yh+gX(Yh).
\]
\begin{enumerate}[start=3]
\item Entonces
\end{enumerate}
\[
[X,gY]h=X(g)Yh+gX(Yh)-gY(Xh)=X(g)Yh+g[X,Y]h.
\]
\begin{enumerate}[start=4]
\item Por lo tanto,
\end{enumerate}
\[
[X,gY]=g[X,Y]+X(g)Y.
\]
\begin{itemize}
\item \textbf{(iv)}
\end{itemize}
\begin{enumerate}
\item Usando \textbf{(iii)} y antisimetría:
\end{enumerate}
\[
[fX,gY]=f[X,gY]-gY(f)X.
\]
\begin{enumerate}[start=2]
\item Luego
\end{enumerate}
\[
[fX,gY]=fg[X,Y]+fX(g)Y-gY(f)X.
\]
\begin{itemize}
\item \textbf{(v)}
\end{itemize}
\begin{enumerate}
\item La prueba de Jacobi es un cálculo. Para $f\in C^\infty(M)$:
\end{enumerate}
\[
\begin{aligned} \big[X,[Y,Z]]f+[Y,[Z,X]]f+[Z,[X,Y]]f&=X[Y,Z]f-[Y,Z]Xf+Y[Z,X]f-[Z,X]Yf+Z[X,Y]f-[X,Y]Zf\\&=XYZf-XZYf-YZXf+ZYXf+YZXf-YXZf\\&\quad -ZXYf+XZYf+ZXYf-ZYXf-XYZf+YXZf\\&=0.\end{aligned}
\]
\bookqed
\end{bookproof}
\begin{bookproposition}{Naturalidad del corchete}
Sean $F:M\to N$ suave, $X,Y\in\mathfrak X(M)$ y $\widetilde X,\widetilde Y\in\mathfrak X(N)$ tales que $X,Y$ estan $F$-relacionados con $\widetilde{X},\widetilde Y$ respectivamente entonces $[X,Y]$ está $F$-relacionado con $[\widetilde X,\widetilde Y]$.

\end{bookproposition}
\begin{bookproof}{}
2. Primero notamos

\[
(dF)_{p}[X,Y]_{p}f=[X,Y]_{p}(f\circ F)=X_{p}(Y(f\circ F))-Y_{p}(X(f\circ F))
\]
\begin{enumerate}[start=3]
\item Tenemos que $(dF)_{p}(X_{p})=\widetilde{X}_{F(p)}$ entonces
\end{enumerate}
\[
\widetilde X_{F(p)}(f)=(\widetilde Xf)(F(p)).
\]
Por lo tanto

\[
X_p(f\circ F)=((\widetilde Xf)\circ F)(p).
\]
Como esto vale para todo $p\in M$, se obtiene la igualdad de funciones

\[
X(f\circ F)=(\widetilde Xf)\circ F.
\]
Analogo con $Y(f\circ F)$. 4. Por lo tanto

\[
\begin{aligned}(dF)_{p}[X,Y]_{p}h=[X,Y]_{p}(h\circ F)&=X_{p}((\widetilde Y f)\circ F)-Y_{p}((\widetilde X f)\circ F)\\&=(dF)_{p}(X_{p})(\widetilde{Y}f)-(dF)_{p}(Y_{p})(\widetilde{X}f)\\&=\widetilde{X}_{F(p)}(\widetilde{Y}f)-\widetilde{Y}_{F(p)}(\widetilde{X}f)\\&=[\widetilde{X},\widetilde{Y}]_{F(p)}f\end{aligned}
\]
mostrando que $[X,Y]$ y $[\widetilde{X},\widetilde{Y}]$ estan $F$-relacionados

\bookqed
\end{bookproof}
\begin{booktheorem}{Derivada de Lie NO FINAL}
Para todo $X,Y\in\mathfrak X(M)$ y para todo $p\in M$, se cumple

\[
[X,Y]_p=(\mathcal L_XY)_p=\lim_{t\to0}\frac{(d\theta_{-t})_{\theta_t(p)}Y_{\theta_t(p)}-Y_p}{t},
\]
donde $\theta$ es el flujo de $X$.

\end{booktheorem}
\begin{bookproof}{}
5. Se calculan ambos lados y se prueba que son iguales. 6. Supongamos primero que $X_p\neq0$. 7. Por el teorema de rectificación de campos, existe una carta suave $(U,\varphi=(x_1,\ldots,x_n))$ centrada en $p$ tal que

\[
X=\frac{\partial}{\partial x_1}.
\]
\begin{enumerate}[start=8]
\item Escribimos
\end{enumerate}
\[
Y=\sum_{j=1}^n a_j\frac{\partial}{\partial x_j},
\]
con $a_1,\ldots,a_n$ funciones suaves sobre $U$. 9. Entonces

\[
[X,Y]=\left[\frac{\partial}{\partial x_1},\sum_{j=1}^n a_j\frac{\partial}{\partial x_j}\right]=\sum_{j=1}^n a_j\left[\frac{\partial}{\partial x_1},\frac{\partial}{\partial x_j}\right]+\sum_{j=1}^n \frac{\partial a_j}{\partial x_1}\frac{\partial}{\partial x_j}.
\]
\begin{enumerate}[start=10]
\item Como los campos coordenados conmutan,
\end{enumerate}
\[
\left[\frac{\partial}{\partial x_1},\frac{\partial}{\partial x_j}\right]=0.
\]
\begin{enumerate}[start=11]
\item Por lo tanto,
\end{enumerate}
\[
[X,Y]_p=\sum_{j=1}^n \frac{\partial a_j}{\partial x_1}(p)\frac{\partial}{\partial x_j}\bigg|_p.
\]
\begin{enumerate}[start=12]
\item Por otro lado, consideremos la curva suave
\end{enumerate}
\[
\alpha(t)=(d\theta_{-t})_{\theta_t(p)}Y_{\theta_t(p)}\in T_pM.
\]
\begin{enumerate}[start=13]
\item Queremos calcular
\end{enumerate}
\[
\alpha'(0)\in T_{\alpha(0)}(T_pM)\simeq T_pM.
\]
\begin{enumerate}[start=14]
\item Entonces
\end{enumerate}
\[
\alpha'(0)=\sum_{j=1}^n c_j\frac{\partial}{\partial x_j}\bigg|_p,
\]
donde

\[
c_j=\alpha'(0)(x_j).
\]
\begin{enumerate}[start=15]
\item Por el lema previo,
\end{enumerate}
\[
\alpha'(0)(x_j)=\frac{d}{dt}\bigg|_{t=0}\alpha(t)(x_j).
\]
\begin{enumerate}[start=16]
\item Luego
\end{enumerate}
\[
\alpha'(0)(x_j)=\frac{d}{dt}\bigg|_{t=0}\left((d\theta_{-t})_{\theta_t(p)}Y_{\theta_t(p)}\right)(x_j).
\]
\begin{enumerate}[start=17]
\item Por definición de diferencial,
\end{enumerate}
\[
\left((d\theta_{-t})_{\theta_t(p)}Y_{\theta_t(p)}\right)(x_j)=Y_{\theta_t(p)}(x_j\circ\theta_{-t}).
\]
\begin{enumerate}[start=18]
\item Como
\end{enumerate}
\[
Y_{\theta_t(p)}=\sum_{i=1}^n a_i(\theta_t(p))\frac{\partial}{\partial x_i}\bigg|_{\theta_t(p)},
\]
tenemos

\[
Y_{\theta_t(p)}(x_j\circ\theta_{-t})=\sum_{i=1}^n a_i(\theta_t(p))\frac{\partial}{\partial x_i}\bigg|_{\theta_t(p)}(x_j\circ\theta_{-t}).
\]
\begin{enumerate}[start=19]
\item En estas coordenadas, como $X=\frac{\partial}{\partial x_1}$, el flujo satisface
\end{enumerate}
\[
x_j\circ\theta_{-t}=x_j-t\delta_{1j}.
\]
\begin{enumerate}[start=20]
\item Por lo tanto,
\end{enumerate}
\[
\frac{\partial}{\partial x_i}\bigg|_{\theta_t(p)}(x_j\circ\theta_{-t})=\delta_{ij}.
\]
\begin{enumerate}[start=21]
\item Así,
\end{enumerate}
\[
Y_{\theta_t(p)}(x_j\circ\theta_{-t})=a_j(\theta_t(p)).
\]
\begin{enumerate}[start=22]
\item Entonces
\end{enumerate}
\[
\alpha'(0)(x_j)=\frac{d}{dt}\bigg|_{t=0}a_j(\theta_t(p))=X_pa_j=\frac{\partial a_j}{\partial x_1}(p).
\]
\begin{enumerate}[start=23]
\item El igual del medio se justifica así: si $v\in T_pM$ y $\gamma$ es una curva con $\gamma(0)=p$ y $\gamma'(0)=v$, entonces
\end{enumerate}
\[
v(f)=\frac{d}{dt}\bigg|_{t=0}f(\gamma(t)).
\]
\begin{enumerate}[start=24]
\item Por lo tanto,
\end{enumerate}
\[
\alpha'(0)=\sum_{j=1}^n \frac{\partial a_j}{\partial x_1}(p)\frac{\partial}{\partial x_j}\bigg|_p=[X,Y]_p.
\]
\begin{enumerate}[start=25]
\item Ahora supongamos que $X_p=0$.
\item En este caso, la curva integral de $X$ que parte de $p$ es constante, luego
\end{enumerate}
\[
\theta_t(p)=p.
\]
\begin{enumerate}[start=27]
\item Entonces la curva queda
\end{enumerate}
\[
\alpha(t)=(d\theta_{-t})_pY_p.
\]
\begin{enumerate}[start=28]
\item Para toda $f\in C^\infty(M)$,
\end{enumerate}
\[
\alpha'(0)(f)=\frac{d}{dt}\bigg|_{t=0}\alpha(t)(f)=\frac{d}{dt}\bigg|_{t=0}Y_p(f\circ\theta_{-t}).
\]
\begin{enumerate}[start=29]
\item Como $Y_p$ es una derivación fija,
\end{enumerate}
\[
\alpha'(0)(f)=Y_p\left(\frac{d}{dt}\bigg|_{t=0}f\circ\theta_{-t}\right).
\]
\begin{enumerate}[start=30]
\item Pero
\end{enumerate}
\[
\frac{d}{dt}\bigg|_{t=0}f(\theta_{-t}(q))=-X_qf.
\]
\begin{enumerate}[start=31]
\item Es decir,
\end{enumerate}
\[
\frac{d}{dt}\bigg|_{t=0}f\circ\theta_{-t}=-Xf.
\]
\begin{enumerate}[start=32]
\item Por lo tanto,
\end{enumerate}
\[
\alpha'(0)(f)=-Y_p(Xf).
\]
\begin{enumerate}[start=33]
\item Como $X_p=0$, también
\end{enumerate}
\[
[X,Y]_pf=X_p(Yf)-Y_p(Xf)=-Y_p(Xf).
\]
\begin{enumerate}[start=34]
\item Así,
\end{enumerate}
\[
\alpha'(0)(f)=[X,Y]_pf.
\]
\begin{enumerate}[start=35]
\item Como esto vale para toda $f\in C^\infty(M)$, concluimos que
\end{enumerate}
\[
\alpha'(0)=[X,Y]_p.
\]
\begin{enumerate}[start=36]
\item Por lo tanto,
\end{enumerate}
\[
(\mathcal L_XY)_p=[X,Y]_p.
\]
\bookqed
\end{bookproof}
\begin{bookproposition}{Criterio de conmutación de flujos (Hulet 9.15)}
Sean $X,Y\in\mathfrak X(M)$ tales que $[X,Y]=0$ y sean $\theta,\phi$ los flujos respectivos. Dado $p\in M$, sean $U$ y $\delta$ los del \hyperref[blk:39cc9a]{Uniformidad conjunta para dos flujos}. Entonces:

\begin{itemize}
\item \textbf{(i)} $(d\theta_s)_qY_q=Y_{\theta_s(q)}$ con $q\in U$ para $s\in(-\delta,\delta)$.
\item \textbf{(ii)} $\theta_t\circ\phi_s=\phi_s\circ\theta_t$ en $U$ para $t,s\in(-\delta,\delta)$.
\end{itemize}
\end{bookproposition}
\begin{bookproof}{}
\begin{itemize}
\item (i)
\end{itemize}
\begin{enumerate}
\item Recordemos la fórmula
\end{enumerate}
\[
[X,Y](p)=\frac{d}{dt}\Big|_{t=0}(d\theta_{-t})_{\theta_t(p)}\big(Y_{\theta_t(p)}\big),\qquad p\in M.
\]
\begin{enumerate}[start=2]
\item Notar que probar $(i)$ equivale a probar que
\end{enumerate}
\[
Y_p=(d\theta_{-s})_{\theta_s(p)}\big(Y_{\theta_s(p)}\big),\qquad \forall s\in (-\delta ,\delta),\forall p\in U
\]
\begin{enumerate}[start=3]
\item En otras palabras, para probar $(i)$ debemos ver que, para cada $p$, la curva
\end{enumerate}
\[
z(s)=(d\theta_{-s})_{\theta_s(p)}\big(Y_{\theta_s(p)}\big)
\]
es constantemente igual a $Y_p$. 4. Calculamos $z'(s)$. Por definición,

\[
z'(s)=\frac{d}{dt}\Big|_{t=0}z(t+s).
\]
\begin{enumerate}[start=5]
\item Ahora,
\end{enumerate}
\[
z(t+s)=(d\theta_{-(t+s)})_{\theta_{t+s}(p)}\big(Y_{\theta_{t+s}(p)}\big).
\]
\begin{enumerate}[start=6]
\item Como
\end{enumerate}
\[
\theta_{t+s}(p)=\theta_t(\theta_s(p))
\]
y

\[
\theta_{-(t+s)}=\theta_{-s}\circ\theta_{-t},
\]
por la regla de la cadena obtenemos

\[
z(t+s)=(d\theta_{-s})_{\theta_s(p)}\left((d\theta_{-t})_{\theta_t(\theta_s(p))}\big(Y_{\theta_t(\theta_s(p))}\big)\right).
\]
\begin{enumerate}[start=7]
\item Por lo tanto,
\end{enumerate}
\[
z'(s)=\frac{d}{dt}\Big|_{t=0}(d\theta_{-s})_{\theta_s(p)}\left((d\theta_{-t})_{\theta_t(\theta_s(p))}\big(Y_{\theta_t(\theta_s(p))}\big)\right).
\]
\begin{enumerate}[start=8]
\item Como $(d\theta_{-s})_{\theta_s(p)}$ es lineal y no depende de $t$,
\end{enumerate}
\[
z'(s)=(d\theta_{-s})_{\theta_s(p)}\left(\frac{d}{dt}\Big|_{t=0}(d\theta_{-t})_{\theta_t(\theta_s(p))}\big(Y_{\theta_t(\theta_s(p))}\big)\right).
\]
\begin{enumerate}[start=9]
\item Por la fórmula del corchete aplicada al punto $\theta_s(p)$,
\end{enumerate}
\[
\frac{d}{dt}\Big|_{t=0}(d\theta_{-t})_{\theta_t(\theta_s(p))}\big(Y_{\theta_t(\theta_s(p))}\big)=[X,Y]_{\theta_s(p)}.
\]
\begin{enumerate}[start=10]
\item Luego
\end{enumerate}
\[
z'(s)=(d\theta_{-s})_{\theta_s(p)}\big([X,Y]_{\theta_s(p)}\big).
\]
\begin{enumerate}[start=11]
\item Como $[X,Y]=0$, concluimos que
\end{enumerate}
\[
z'(s)=0,\qquad \forall s\in (-\delta ,\delta )
\]
osea $z(s)$ es constante 12. Además,

\[
z(0)=(d\theta_0)_p(Y_p)=Y_p.
\]
\begin{enumerate}[start=13]
\item Por lo tanto,
\end{enumerate}
\[
z(s)=Y_p,\qquad \forall s,
\]
es decir,

\[
(d\theta_{-t})_{\theta_t(p)}\big(Y_{\theta_t(p)}\big)=Y_p.
\]
\begin{enumerate}[start=14]
\item Aplicando $(d\theta_t)_p$ a ambos lados, obtenemos
\end{enumerate}
\[
(d\theta_t)_p(Y_p)=Y_{\theta_t(p)}.
\]
\begin{enumerate}[start=15]
\item \textbf{(ii)} Ahora fijemos $t$ y $q$. Consideramos
\end{enumerate}
\[
\alpha(s)=\theta_t(\phi_s(q)),
\]
\[
\beta(s)=\phi_s(\theta_t(q)).
\]
\begin{enumerate}[start=16]
\item Tenemos
\end{enumerate}
\[
\alpha(0)=\theta_t(q)=\beta(0).
\]
\begin{enumerate}[start=17]
\item Además,
\end{enumerate}
\[
\alpha'(s)=(d\theta_t)_{\phi_s(q)}Y_{\phi_s(q)}=Y_{\theta_t(\phi_s(q))}=Y_{\alpha(s)},
\]
donde usamos \textbf{(i)}. 18. Por otro lado,

\[
\beta'(s)=Y_{\phi_s(\theta_t(q))}=Y_{\beta(s)}.
\]
\begin{enumerate}[start=19]
\item Entonces $\alpha$ y $\beta$ son curvas integrales de $Y$ con el mismo dato inicial.
\item Por unicidad,
\end{enumerate}
\[
\theta_t(\phi_s(q))=\phi_s(\theta_t(q)).
\]
\bookqed
\end{bookproof}
\begin{bookproposition}{(Hulet 9.16)}
Sean $X,Y\in\mathfrak{X}(M)$, cuyos flujos conmutan. Entonces

\[
[X,Y]=0.
\]
\end{bookproposition}
\begin{bookproof}{}
\begin{enumerate}
\item Sean $\theta$ y $\phi$ los flujos de $X$ e $Y$, respectivamente. Fijemos $p\in M$. Por el lema de uniformidad conjunta de los flujos, existen un abierto $U$ que contiene a $p$ y $\delta>0$ tales que, para $|s|,|t|<\delta$, las composiciones $\theta_t\circ\phi_s$ y $\phi_s\circ\theta_t$ están definidas en $U$.
\item Como los flujos conmutan, para $|s|,|t|<\delta$ se cumple
\end{enumerate}
\[
\theta_t(\phi_s(p))=\phi_s(\theta_t(p)).
\]
\begin{enumerate}[start=3]
\item Derivando respecto de $s$ en $s=0$, obtenemos
\end{enumerate}
\[
(d\theta_t)_p(Y_p)=(d\theta_t)_{\phi_0(p)}\left(\frac{d}{ds}\bigg|_{s=0}\phi_s(p)\right)=\frac{d}{ds}\bigg|_{s=0}\theta_t(\phi_s(p))=\frac{d}{ds}\bigg|_{s=0}\phi_s(\theta_t(p))=Y_{\theta_t(p)}.
\]
\begin{enumerate}[start=4]
\item Aplicamos $(d\theta_{-t})_{\theta_t(p)}$ a ambos lados. Entonces
\end{enumerate}
\[
\begin{aligned}(d\theta_{-t})_{\theta_t(p)}\left((d\theta_t)_p(Y_p)\right)&=(d\theta_{-t})_{\theta_t(p)}\left(Y_{\theta_t(p)}\right),\\d(\theta_{-t}\circ\theta_t)_p(Y_p)&=(d\theta_{-t})_{\theta_t(p)}\left(Y_{\theta_t(p)}\right),\\Y_p&=(d\theta_{-t})_{\theta_t(p)}\left(Y_{\theta_t(p)}\right).\end{aligned}
\]
\begin{enumerate}[start=5]
\item Por lo tanto, la curva
\end{enumerate}
\[
t\mapsto(d\theta_{-t})_{\theta_t(p)}\left(Y_{\theta_t(p)}\right)
\]
es constante para $|t|<\delta$. 6. Usando la fórmula del corchete en términos del flujo de $X$,

\[
[X,Y]_p=\frac{d}{dt}\bigg|_{t=0}(d\theta_{-t})_{\theta_t(p)}\left(Y_{\theta_t(p)}\right)=0.
\]
\begin{enumerate}[start=7]
\item Como $p\in M$ era arbitrario, concluimos que
\end{enumerate}
\[
[X,Y]=0.
\]
\bookqed
\end{bookproof}
\begin{bookproposition}{Rectificación simultánea (Hulet 9.18)}
Sean $X_1,\ldots,X_k\in\mathfrak X(M)$ linealmente independientes en un abierto de $M$ en $p$ tales que

\[
[X_{i},X_{j}]=0\quad 1\leq i,j\leq k
\]
Entonces existe un sistema coordenado $(U,\varphi=(x_1,\ldots,x_n))$ centrado en $p$ tal que

\[
X_1|_U=\frac{\partial}{\partial x_1},\ldots,\qquad X_k|_U=\frac{\partial}{\partial x_k}
\]
\end{bookproposition}
\begin{bookproof}{}
8. Por \hyperref[blk:39cc9a]{Uniformidad conjunta para dos flujos}. Existe $\varepsilon>0$ y un abierto $V$ de $p$ tal que la composición

\[
\theta^{1}_{t_1}\circ\cdots\circ\theta^k_{t_k}
\]
está bien definida sobre $V$ para todos $t_1,\ldots,t_k\in(-\varepsilon,\varepsilon)$. 9. Como los campos son linealmente independientes. Achicando $V$ si hace falta, tomamos una carta cúbica centrada en $p$,

\[
(V,\psi=(x_1,\ldots,x_n)),
\]
tal que

\[
\left.\frac{\partial}{\partial x_i}\right|_p=X_i|_p,\qquad i=1,\ldots,k.
\]
por \hyperref[blk:047506]{Lema} 10. Nos enfocamos en la rebanada

\[
R=\{(r_1,\ldots,r_n)\in C_{\epsilon}^{n} (0):r_1=\cdots=r_k=0\}\simeq C_{\epsilon}^{n-k} (0)\subseteq \psi(V)
\]
que tiene dimensión $n-k$. 11. Armamos el mapa

\[
\sigma:(-\varepsilon,\varepsilon)^k\times C_{\epsilon}^{n-k}\subseteq \mathbb{R}^{n}  \to M
\]
por

\[
\sigma(t_1,\ldots,t_k,r^{k+1},\ldots,r^{n}  )=\theta^1_{t_1}\circ\cdots\circ\theta^k_{t_k}(\psi^{-1}(0,\ldots 0,r^{k+1},\ldots,r^{n}  )).
\]
Notar que esto este bien definido por que al ser carta cubica $\psi ^{-1}(0,\ldots,0,r^{k+1},\ldots r^{n})$ cae en $V$ 12. Como en la vez anterior, vamos a ver que $\sigma$ es un difeomorfismo entre un abierto de $\mathbb{R}^n$ y un abierto de $M$ que contiene a $p$, mostrando que

\[
(d\sigma)_0\left(\left.\frac{\partial}{\partial r^i}\right|_0\right)=\begin{cases} X_i|_p, & i=1,\ldots,k,\\[4pt] \left.\dfrac{\partial}{\partial x^i}\right|_p, & i>k.\end{cases}
\]
\begin{enumerate}[start=13]
\item En efecto, si $i\leq k$,
\end{enumerate}
\[
\begin{aligned}(d\sigma)_0\left(\left.\frac{\partial}{\partial t^i}\right|_0\right)&=\left.\frac{d}{ds}\right|_{s=0}\sigma(se_i)\\&=\left.\frac{d}{ds}\right|_{s=0}\sigma(0,\ldots,0,s,0,\ldots,0)\\&=\left.\frac{d}{ds}\right|_{s=0}\left(\theta^{1}_{0}\circ\cdots\circ\theta^{i}_{s}\circ\cdots\circ\theta^{k}_{0}\right)\left(\psi^{-1}(0,\ldots,0)\right)\\&=\left.\frac{d}{ds}\right|_{s=0}\theta^{i}_{s}(p)\\&=X_i|_p\end{aligned}
\]
\begin{enumerate}[start=14]
\item Si $i>k$,
\end{enumerate}
\[
\begin{aligned}(d\sigma)_0\left(\left.\frac{\partial}{\partial r^i}\right|_0\right)&=\left.\frac{d}{ds}\right|_{s=0}\sigma(se_i)\\&=\left.\frac{d}{ds}\right|_{s=0}\sigma(0,\ldots,0,s,0,\ldots,0)\\&=\left.\frac{d}{ds}\right|_{s=0}(\theta^{1}_{0}\circ\cdots\circ\theta^{k}_{0})(\psi^{-1}(0,\ldots,0,s,0,\ldots,0)\\&=\left.\frac{d}{ds}\right|_{s=0}\psi^{-1}(0,\ldots,0,s,0,\ldots,0)\\&=(d\psi^{-1})_0\left(\left.\frac{\partial}{\partial r^i}\right|_0\right)\\&=\left.\frac{\partial}{\partial x^i}\right|_p\end{aligned}
\]
este ultimo igual sale rapido escribiendo a $(d\psi^{-1})_0\left(\left.\frac{\partial}{\partial r^i}\right|_0\right)$ en coordenadas del $T_{\phi ^{-1}(0)}M$ 15. Ahora notamos que por paso 2.

\[
\left\{ X_{1}|_{p},\ldots,X_{k}|_{p},\frac{\partial}{\partial x^{k+1}}\bigg|_{p},\ldots, \frac{\partial}{\partial x^{n}}|_{p}\right\}=\left\{ \frac{\partial}{\partial x^{n}}\bigg|_{p},\ldots, \frac{\partial}{\partial x^{k}}\bigg|_{p} , \frac{\partial}{\partial x^{k+1}}|_{p}\ldots\frac{\partial}{\partial x^{n}}\bigg|_{p} \right\}
\]
\begin{enumerate}[start=16]
\item Entonces $(d\sigma)_{0}$ es isomorfismo lineal por que manda base $\{ \frac{\partial}{\partial r_{i}} \}_{i<n}$ (notar aca llamo $t_{i}:=r^{i}$ es solo notacion para cordenadas en $\mathbb{R}$) en una base $\{ \frac{\partial}{\partial x^{n}}|_{p},\ldots, \frac{\partial}{\partial x^{k}}|_{p} , \frac{\partial}{\partial x^{k+1}}|_{p}\ldots\frac{\partial}{\partial x^{n}}|_{p} \}$
\item Estamos entonces en condiciones de aplicar el Teorema de la Función Inversa. Existe un abierto $\widetilde U\subseteq\mathbb{R}^n$ que contiene a $0$ y un abierto $U\subseteq M$ que contiene a $p$ tales que
\end{enumerate}
\[
\sigma|_{\widetilde U}:\widetilde U\longrightarrow U
\]
es un difeomorfismo. 18. Como $\sigma$ es un difeomorfismo entre $\widetilde U$ y $U$, llamamos $\sigma^{-1}=\varphi$, de modo que $(U,\varphi=y^{1},\ldots,y^{i})$ es una carta coordenada. 19. Recordamos nuestro objetivo es probar que

\[
X_i|_U=\frac{\partial}{\partial y^i},\qquad i=1,\ldots,k.
\]
\begin{enumerate}[start=20]
\item Equivalentemente, queremos ver que para todo $(r^1,\ldots,r^n)\in\widetilde U$, e $1\leq i\leq k$
\end{enumerate}
\[
(d\sigma)_{(r^1,\ldots,r^n)}\left(\left.\frac{\partial}{\partial r^i}\right|_{(r^1,\ldots,r^n)}\right)=(X_i)_{\sigma(r^1,\ldots,r^n)}.
\]
por que si esto es cierto usamos que

\[
(d\sigma)_{(r^1,\ldots,r^n)}\left(\left.\frac{\partial}{\partial r^i}\right|_{(r^1,\ldots,r^n)}\right)=(d\varphi ^{-1})_{(r^1,\ldots,r^n)}\left(\left.\frac{\partial}{\partial r^i}\right|_{(r^1,\ldots,r^n)}\right)=\frac{\partial}{\partial y^{i} }\bigg|_{\varphi ^{-1}(r)=\sigma(r)}
\]
\begin{enumerate}[start=21]
\item Primero notemos que, como
\end{enumerate}
\[
[X_i,X_j]=0
\]
los flujos conmutan localmente por \hyperref[blk:26db5d]{Criterio de conmutación de flujos (Hulet 9.15)} 22. Fijemos $1\leq i\leq k$. Entonces

\[
\begin{aligned}(d\sigma)_{(r^1,\ldots,r^n)}\left(\left.\frac{\partial}{\partial r^i}\right|_{(r^1,\ldots,r^n)}\right)&=\left.\frac{d}{ds}\right|_{s=0}\sigma(r^1,\ldots,r^i+s,\ldots,r^n)\\&=\left.\frac{d}{ds}\right|_{s=0}\left(\theta^{1} _{r^1}\circ\cdots\circ\theta^{i} _{r^i+s}\circ\cdots\circ\theta_{k,r^k}\right)\left(\psi^{-1}(0,\ldots,0,r^{k+1},\ldots,r^n)\right)\\&=\left.\frac{d}{ds}\right|_{s=0}\left(\theta^{1} _{r^1}\circ\cdots\circ\theta^{i} _{r^i}\circ\theta^{i} _{s}\circ\cdots\circ\theta^{k} _{r^k}\right)\left(\psi^{-1}(0,\ldots,0,r^{k+1},\ldots,r^n)\right)\\&=\left.\frac{d}{ds}\right|_{s=0}\theta^{i} _{s}\left(\sigma(r^1,\ldots,r^n)\right)\\&=(X_i)_{\sigma(r^1,\ldots,r^n)}.\end{aligned}
\]
donde usamos que los flujos conmutan

\bookqed
\end{bookproof}
\section{Conjunto de formas diferenciales}
\begin{booktheorem}{Suavidad de $k$-formas}
Sea $\omega:M\to\Lambda^k(T^*M)$ una $k$-forma. Las siguientes afirmaciones son equivalentes:

\begin{itemize}
\item (a) $\omega$ es suave.
\item (b) Para toda carta $(U,\varphi=(x_1,\ldots,x_n))$, las coordenadas de $\omega$ con respecto al marco coordenado $\{dx_{i_1}\wedge\cdots\wedge dx_{i_k}\}_{1\le i_1<\cdots<i_k\le n}$ son funciones suaves.
\item (c) Para cualquier $k$ campos suaves $X_1,\ldots,X_k\in\mathfrak X(M)$, la funcion $\omega(X_1,\ldots,X_k)\in C^\infty(M)$.
\end{itemize}
\end{booktheorem}
\begin{bookproof}{}
\begin{itemize}
\item $(a\Rightarrow b)$.
\end{itemize}
\begin{enumerate}
\item Supongamos que $\omega:M\to\Lambda^k(T^*M)$ es suave.
\item Sea $(U,\varphi=(x^1,\ldots,x^n))$ una carta de $M$.
\item Esta carta induce una carta $(\widetilde U,\widetilde\varphi)$ de $\Lambda^k(T^*M)$, donde $\widetilde U=\pi^{-1}(U)$. Usamos la notación de multiíndices $I=(i_1,\ldots,i_k)$, con $i_1<\cdots<i_k$, y escribimos $dx^I=dx^{i_1}\wedge\cdots\wedge dx^{i_k}$. La carta inducida está dada por
\end{enumerate}
\[
\widetilde\varphi\left(\sum_I C_I\,dx^I|_p\right)=\left(x^1(p),\ldots,x^n(p),(C_I)_I\right).
\]
\begin{enumerate}[start=4]
\item Como $\omega$ es suave, su restricción $\omega|_U:U\to\widetilde U$ es suave. Por lo tanto, como $\widetilde\varphi$ es suave, la composición
\end{enumerate}
\[
\widetilde\varphi\circ\omega|_U:U\to\mathbb R^{n+\binom nk}
\]
es suave. 5. Escribimos, para cada $p\in U$,

\[
\omega_p=\sum_I C_I(p)\,dx^I|_p.
\]
\begin{enumerate}[start=6]
\item Entonces
\end{enumerate}
\[
(\widetilde\varphi\circ\omega|_U)(p)=\left(x^1(p),\ldots,x^n(p),(C_I(p))_I\right).
\]
\begin{enumerate}[start=7]
\item Como $\widetilde\varphi\circ\omega|_U$ es suave, todas sus funciones coordenadas son suaves. En particular, para cada multiíndice $I=(i_1,\ldots,i_k)$ con $1\le i_1<\cdots<i_k\le n$, la función $C_I:U\to\mathbb R$ es suave.
\item Como la carta $(U,\varphi)$ era arbitraria, las coordenadas $C_I$ de $\omega$ respecto de todo marco coordenado $\{dx^I\}_I$ son funciones suaves.
\end{enumerate}
\begin{itemize}
\item $(b\Rightarrow a)$.
\end{itemize}
\begin{enumerate}
\item Supongamos que, para toda carta $(U,\varphi=(x^1,\ldots,x^n))$, las coordenadas $C_I:U\to\mathbb R$ de $\omega$ son funciones suaves.
\item Fijemos una de estas cartas y consideremos la carta inducida $(\widetilde U,\widetilde\varphi)$ de $\Lambda^k(T^*M)$, con $\widetilde U=\pi^{-1}(U)$.
\item Para cada $p\in U$, escribimos
\end{enumerate}
\[
\omega_p=\sum_I C_I(p)\,dx^I|_p.
\]
\begin{enumerate}[start=4]
\item Por la definición de la carta inducida,
\end{enumerate}
\[
(\widetilde\varphi\circ\omega|_U)(p)=\left(x^1(p),\ldots,x^n(p),(C_I(p))_I\right).
\]
\begin{enumerate}[start=5]
\item Las funciones $x^1,\ldots,x^n$ son suaves y, por hipótesis, todas las funciones $C_I$ son suaves. Por lo tanto, $\widetilde\varphi\circ\omega|_U$ es suave.
\item Como $\widetilde\varphi$ es una carta, concluimos que $\omega|_U:U\to\widetilde U$ es suave.
\item Como esto vale para toda carta $(U,\varphi)$ de $M$, concluimos que $\omega:M\to\Lambda^k(T^*M)$ es suave.
\end{enumerate}
\begin{itemize}
\item $(b\Rightarrow c)$
\end{itemize}
\begin{enumerate}
\item Tomamos un $p\in M$ ahora tenemos $(U,\phi=x_{1},\ldots,x_{n})$ carta entonces por (b)
\end{enumerate}
\[
\omega|_U=\sum_{i_1<\cdots<i_k}C_{i_1\cdots i_k}\,dx_{i_1}\wedge\cdots\wedge dx_{i_k},\qquad C_{i_1\cdots i_k}\in C^\infty(U)
\]
\begin{enumerate}[start=2]
\item Ademas
\end{enumerate}
\[
X_{i}=\sum^{n}_{j=1}a_{ij}\frac{\partial}{\partial x_{j}}
\]
\begin{enumerate}[start=3]
\item Luego
\end{enumerate}
\[
\omega|_{U}(X_{1},\ldots X_{k})=\sum_{i_1<\cdots<i_k}C_{i_1\cdots i_k}\,dx_{i_1}\wedge\cdots\wedge dx_{i_k}\left(\sum^{n}_{j=1}a_{1j}\frac{\partial}{\partial x_{j}},\ldots,\sum^{n}_{j=1}a_{kj}\frac{\partial}{\partial x_{j}}\right)
\]
\begin{enumerate}[start=4]
\item Ahora notamos que
\end{enumerate}
\[
\begin{aligned} dx_{i_1}\wedge\cdots\wedge dx_{i_k}\left(\sum^{n}_{j=1}a_{1j}\frac{\partial}{\partial x_{j}},\ldots,\sum^{n}_{j=1}a_{kj}\frac{\partial}{\partial x_{j}}\right)&=\det\left( dx_{i_{r}}\left( \sum^{n}_{j=1}a_{sj}\frac{\partial}{\partial x_{j}} \right) \right)_{r,s=1}^{k} \\ & =\det\left( \sum^{n}_{j=1} a_{sj}dx_{i_{r}}\left(\frac{\partial}{\partial x_{j}}\right) \right)_{r,s=1}^{k}\\&=\det(a_{s i_{r}})_{r,s=1}^{k} \end{aligned}
\]
\begin{enumerate}[start=5]
\item Pero entonces
\end{enumerate}
\[
\omega|_{U}(X_{1},\ldots,X_{k})(p)=(\omega|_{U})_{p}(X_{1},\ldots,X_{k})=\sum_{i_{1}<\ldots<i_{k}} C_{i_{1}\ldots i_{k}}(p)\det(a_{s i_{r}}(p))
\]
\begin{enumerate}[start=6]
\item Pero $C_{i_{1}\ldots i_{k}}$ es suave como funcion de $p$ por hipotesis, lo mismo para $\det(a_{s i_{r}})_{r,s=1}^{k}$ como funcion de $p$ por que $X_{i}$ son campos suaves. Y $\det$ es un polinomio de sus entradas.
\item Entonces $\omega|_{U}(X_{1},\ldots,X_{k})(p)$ es suave y obviamente como esto lo podemos hacer para cualquier $p$ tenemosq ue
\end{enumerate}
\[
\omega(X_{1},\ldots,X_{k})(p)
\]
es suave

\begin{itemize}
\item $(c\Rightarrow b)$
\end{itemize}
\begin{enumerate}
\item Tomamos una carta $(U,x_{1},\ldots,x_{n})$ como $B_k=\{dx_{i_1}|_p\wedge\cdots\wedge dx_{i_k}|_p:1\le i_1<\cdots<i_k\le n\}$ es base de $\Lambda^{k}(T_{p}^{*}M)$ tenemos
\end{enumerate}
\[
\omega|_{U}=\sum_{i_{1}<\ldots<i_{k}} C_{i_{1}\ldots i_{k}}\ .dx_{i_{1}}\wedge\ldots\wedge dx_{i_{k}}
\]
con $C_{i_{1}\ldots i_{k}}:U \rightarrow\mathbb{R}$. Bastaria mostrar que son suaves 2. Ahora tenemos que

\[
\begin{aligned}\omega|_{U}\left(\frac{\partial}{\partial x_{j_{1}}},\ldots, \frac{\partial}{\partial x_{j_{k}}}\right)(p)& =\omega_{p}\left(\frac{\partial}{\partial x_{j_{1}}}\bigg|_{p},\ldots, \frac{\partial}{\partial x_{j_{k}}}\bigg|_{p}\right)\\&=\sum^{\infty}_{i_{1}<\ldots<i_{k}}C_{i_{1}\ldots i_{k}}(p) \ .dx_{i_{1}}\wedge\ldots\wedge dx_{i_{k}}\left(\frac{\partial}{\partial x_{j_{1}}}\bigg|_{p},\ldots, \frac{\partial}{\partial x_{j_{k}}}\bigg|_{p}\right)\\&=C_{I}(p)\end{aligned}
\]
con $I=J$ osea el unico caso donde el determinante no es $0$ que es cuando coinciden los multi indices, $j_{1}=i_{1}\ldots j_{k}=i_{k}$\\
3. Ahora lo unico que tenemos que corregir es que $\frac{\partial}{\partial x_{i}}$ son campos suaves sobre $U$  y para usar la hipotesis las necesito suaves sobre $M$ 4. Tomamos $q\in U$ y ahora tenemos $V_{q}\subseteq U$ abierto de $q$ y $f\in C^{\infty}(M)$ tal que $f|_{\overline{V_{q}}}\equiv 1$ y $\operatorname{supp}f\subseteq U$ y defino

\[
X_{i}(p)=\begin{cases} f(p)\frac{\partial}{\partial x_{i}}\bigg|_{p} & p\in U \\0 & x\not\in U \end{cases}
\]
que sabemos es un campo suave sobre todo $M$ y cerca de $p$ (osea en algun abierto dentro de $\overline{V_{q}}$, por ejemplo $V_{q}$ ) vale exactamente $\frac{\partial}{\partial x_{i}}\bigg|_{q}$ 5. Luego haciendo la misma cuenta que en 2. pero con $X_{i}$ tenemos que

\[
\sum_{i_{1}<\ldots<i_{k}} C_{i_{1}\ldots i_{k}}(q)
\]
es suave para todo $q\in V_{p}$ 6. Y como esto lo podemos hacer para cualquier $q\in U$ es suave en todo $U$ como queriamos

\bookqed
\end{bookproof}
\begin{bookproposition}{Ejercicio del practico}
Probar que si $\omega\in\Omega^k(M)$ y $\eta\in\Omega^\ell(M)$, entonces $\omega\wedge\eta\in\Omega^{k+\ell}(M)$.

\end{bookproposition}
\begin{bookproof}{}
\begin{enumerate}
\item En coordenadas, escribimos
\end{enumerate}
\[
\omega=\sum_I a_I\,dx^I,\qquad \eta=\sum_J b_J\,dx^J,
\]
con $a_I,b_J$ funciones suaves. 2. Entonces

\[
\omega\wedge\eta=\sum_{I,J}a_Ib_J\,dx^I\wedge dx^J.
\]
\begin{enumerate}[start=3]
\item Los términos que tienen diferenciales repetidos dan $0$, y los demás se reordenan posiblemente cambiando el signo. Pero los nuevos coeficientes siguen siendo productos $\pm a_Ib_J$ de funciones suaves.
\item Por lo tanto, en cada carta $\omega\wedge\eta$ tiene coordenadas suaves. Luego $\omega\wedge\eta\in\Omega^{k+\ell}(M)$.
\end{enumerate}
\bookqed
\end{bookproof}
\begin{booklemma}{La diferencial exterior es local, lema previo}
Sea $d:\Omega^*(M)\to\Omega^*(M)$ satisfaciendo (i)-(iv). Sean $\omega_1,\omega_2\in\Omega^*(M)$ tales que coinciden en un abierto $U$ de $M$. Entonces

\[
(d\omega_1)|_U=(d\omega_2)|_U.
\]
\end{booklemma}
\begin{bookproof}{}
\begin{enumerate}
\item Como $d$ es transformacion lineal, podemos mirar $\omega=\omega_1-\omega_2$ y asi es suficiente con ver que si $\omega|_U=0$ (Osea $\omega_{1}$ y $\omega_{2}$ coinciden en $U$ abierto), entonces $(d\omega)|_U=0$. Con lo cual se cumpliria la igualdad pedida
\item Aca hacemos un truco clasico de tensores y es volver todo cero usando una funcion de levantamiento.
\item Sabemos que las funciones de levantamiento preservan en un abierto contenido en otro y matan por "afuera", pero aca $\omega$ ya es cero en $U$. La idea es "volverlo cero" en todo $M$.
\item Sea $p\in U$ arbitrario. Sea $V$ abierto de $M$ con $p\in V$ y $\overline V\subset U$. Sea $f\in C^\infty(M)$ funcion de levantamiento tal que $f|_{\overline V}\equiv1$ y $\operatorname{supp}f\subset U$.
\item Como $f$ es suave, entonces $f\omega\in\Omega^*(M)$, ademas
\end{enumerate}
\[
(f\omega)_q=f(q)\omega_q=\begin{cases}f(q)0,&q\in U,\\0\omega_q,&q\notin U,\end{cases}
\]
porque $\operatorname{supp}f\subset U$. Asi $f\omega=0$ en todo $M$. 6. Como $d$ es lineal, $d(f\omega)=d(0)=0$. Pero por otro lado

\[
d(f\omega)=d(f\wedge\omega)=(df)\wedge\omega+(-1)^{0.k} f\,d\omega.
\]
\begin{enumerate}[start=7]
\item Luego
\end{enumerate}
\[
0=(df)\wedge\omega+f\,d\omega.
\]
\begin{enumerate}[start=8]
\item Evaluando en $p\in U$, tenemos
\end{enumerate}
\[
0=((df)\wedge\omega)_p+f(p)(d\omega)_p=(df)_p\wedge\omega_p+(d\omega)_p=(d\omega)_p,
\]
porque $\omega_p=0$ y $f(p)=1$. Como $p$ era arbitrario, $(d\omega)|_U=0$.

\bookqed
\end{bookproof}
\begin{booktheorem}{Derivada exterior}
Sea $M$ variedad suave. Entonces existe una unica transformacion lineal

\[
d:\Omega^*(M)\to\Omega^*(M)
\]
llamada derivada exterior, que satisface:

\begin{enumerate}
\item $d:\Omega^k(M)\to\Omega^{k+1}(M)$.
\item $d(f)$ es $df$, diferencial de $f$, para toda $f\in\Omega^0(M)=C^\infty(M)$.
\item Si $\omega\in\Omega^k(M)$ y $\theta\in\Omega(M)$, entonces
\end{enumerate}
\[
d(\omega\wedge\theta)=(d\omega)\wedge\theta+(-1)^k\omega\wedge d\theta.
\]
\begin{enumerate}[start=4]
\item $d^2=0$.
\end{enumerate}
\end{booktheorem}
\begin{bookproof}{}
\begin{itemize}
\item \textbf{Unicidad}
\end{itemize}
\begin{enumerate}
\item La idea de la prueba es construir $d$ de forma local, usando las condiciones que debemos satisfacer: (ii)-(iv). Hacemos ingenieria inversa para saber como deberia ser $d$, localmente. Primero, notemos:
\item Asumiendo existencia veamos primero la unicidad, que nos va dar pistas de como definir $d$.
\item Como $d$ se va definir localmente , fijemos una carta suave $(U,\varphi=(x_1,\ldots,x_n))$ y tomemos $\omega\in\Omega^k(M)$ arbitrario y llegemos a una formula para $(d\omega)|_U$.
\item Tenemos
\end{enumerate}
\[
\omega|_U=\sum_{1\le i_1<\cdots<i_k\le n}C_{i_1\ldots i_k}\,dx_{i_1}\wedge\cdots\wedge dx_{i_k}.
\]
\begin{enumerate}[start=5]
\item Pero $d$ actua sobre cosas definidas sobre todo $M$, entonces no podemos aplicar directamente $d$ a la expresion anterior. Para hacerlo bien, debemos extender dicha expresion a todo $M$, usando una funcion de levantamiento.
\item Sea $V$ abierto de $M$ con $\overline V\subset U$, y sea $f\in C^\infty(M)$ tal que $f|_{\overline V}\equiv1$ y $\operatorname{supp}f\subset U$. Extendemos las $C_{i_1\ldots i_k}$ y $x_1,\ldots,x_n$ definiendo:
\end{enumerate}
\[
f\cdot x_i=\begin{cases}f(q)x_i(q),&q\in U,\\0,&\text{en otro caso},\end{cases}\qquad f\cdot C_{i_1\ldots i_k}=\begin{cases}f(q)C_{i_1\ldots i_k}(q),&q\in U,\\0,&\text{en otro caso}.\end{cases}
\]
\begin{enumerate}[start=7]
\item Estas nuevas funciones son suaves (hemos hecho cuentas similares en otras ocaciones)
\item Consideremos la nueva $k$-forma:
\end{enumerate}
\[
\widetilde{\omega}=\sum_{1\le i_1<\cdots<i_k\le n}\underbrace{f.C_{i_1\cdots i_k}}_{\text{smooth function}}\ \ \underbrace{\underbrace{d(f.x_{i_1})}_{\text{smooth }1\text{-form}}\wedge\cdots\wedge\underbrace{d(f.x_{i_k})}_{\text{smooth }1\text{-form}}}_{\text{smooth }k\text{-form}}.
\]
$d(f.x_{i_{j}})$ son $1$-formas por lo mismo que \hyperref[blk:171d2a]{Motivacion} y luego usamos \hyperref[blk:ce25ad]{Ejercicio del practico} para el wedge de todas esas $1$-formas y llegamos a la $k$-forma\\
9. Notar que $\widetilde\omega$ coincide con $\omega$ en el abierto $V$. En efecto, si $q\in V$, tenemos $f(q)=1$ y asi

\[
(f\cdot C_{i_1\ldots i_k})(q)=f(q).C_{i_1\ldots i_k}(q)=C_{i_1\ldots i_k}(q)
\]
\begin{enumerate}[start=10]
\item Y ademas $d(f\cdot x_i)_q=(dx_i)_q$ por que
\end{enumerate}
\[
d(f.x_{i})_{q}\left(\frac{\partial}{\partial x_{i}}\bigg|_{q}\right)=\left(\frac{\partial}{\partial x_{i}}\bigg|_{q}\right)(f.x_{i})=\left(\frac{\partial}{\partial x_{i}}\bigg|_{q}\right)(x_{i})=(dx_{i})_{q}\left(\frac{\partial}{\partial x_{i}}\bigg|_{q}\right)
\]
usando definicion de diferencial y en el segundo igual como $f.x_{i}\equiv x_{i}$ en el abierto $V$ sabemos que $\left(\frac{\partial}{\partial x_{i}}\bigg|_{q}\right)$ actua de la misma forma en ambos (germen local) 11. Pero entonces

\[
(d(f.x_{i_{1}})\wedge\ldots\wedge d(f.x_{i_{k}}))_{q}=(d(x_{i_{1}})\wedge\ldots\wedge d(x_{i_{k}}))_{q}
\]
usando la definicion de producto cuña de $k$-formas y la asociatividad 12. Luego $w$ y $\widetilde{w}$  coinciden en un abierto $V$. Por tanto, por \hyperref[blk:ef3cfa]{La diferencial exterior es local, lema previo}

\[
(d\omega)|_V=(d\widetilde\omega)|_V.
\]
\begin{enumerate}[start=13]
\item Lo bueno es que la expresion para $\widetilde\omega$ esta definida sobre todo $M$. Por lo tanto podemos aplicar $d$ y obtenemos:
\end{enumerate}
\[
(d\widetilde\omega)=d\sum fC_{i_1\ldots i_k}\,d(fx_{i_1})\wedge\cdots\wedge d(fx_{i_k})=d\sum fC_{i_1\ldots i_k}\wedge d(fx_{i_1})\wedge\cdots\wedge d(fx_{i_k})
\]
el ultimo igual por que $fC_{i_1\ldots i_k}$ es una $0$-forma 14. Por linealidad de $d$ y usando la propiedad $3$ de esta misma afirmacion y considerando que $fC_{i_1\ldots i_k}$ es una $0$-forma, queda

\[
(d\widetilde\omega)=\sum d(fC_{i_1\ldots i_k})\wedge d(fx_{i_1})\wedge\cdots\wedge d(fx_{i_k})+\sum fC_{i_1\ldots i_k}\wedge d( d(fx_{i_1})\wedge\cdots\wedge d(fx_{i_k}))
\]
\begin{enumerate}[start=15]
\item Ahora como $d^{2}=0$
\end{enumerate}
\[
(d\widetilde\omega)=\sum d(fC_{i_1\ldots i_k})\wedge d(fx_{i_1})\wedge\cdots\wedge d(fx_{i_k})
\]
\begin{enumerate}[start=16]
\item Finalmente como por paso 12. $(d\omega)|_{V}=(d \widetilde{\omega})|_{V}$ y por paso 16. y por la misma idea que en 11. tenemos
\end{enumerate}
\[
(d\omega)|_V=(d \widetilde{\omega})|_{V}=\sum d(C_{i_1\ldots i_k})\wedge dx_{i_1}\wedge\cdots\wedge dx_{i_k}.
\]
\begin{enumerate}[start=17]
\item De lo que resulta que $d$ es unico pues la formula anterior solo depende de la diferencial usual de funciones suaves.
\item Notar que la formula anterior da pistas de como se debe definir la transformacion lineal $d:\Omega^*(M)\to\Omega^*(M)$.
\end{enumerate}
\begin{itemize}
\item \textbf{Existencia}
\end{itemize}
\begin{enumerate}
\item Construyamos $d$ localmente y luego lo hacemos de forma global. Fijemos una carta $(U,\varphi=(x_1,\ldots,x_n))$.
\item Dada una $k$-forma $\omega\in\Omega^k(U)$,
\end{enumerate}
\[
\omega=\sum C_{i_1\ldots i_k}\,dx_{i_1}\wedge\cdots\wedge dx_{i_k},\qquad C_{i_1\ldots i_k}\in C^\infty(U),
\]
\begin{enumerate}[start=3]
\item Definimos $d_U$ por
\end{enumerate}
\[
d_U\omega:=\sum dC_{i_1\ldots i_k}\wedge dx_{i_1}\wedge\cdots\wedge dx_{i_k},
\]
donde $dC_{i_1\ldots i_k}$ es la diferencial de $C_{i_1\ldots i_k}$, por lo tanto una $1$-forma y luego su base es $\{ dx_{1},\ldots,dx_{n} \}$

\[
dC_{i_1\ldots i_k}=\sum_{j=1}^n\frac{\partial C_{i_1\ldots i_k}}{\partial x_j}\,dx_j.
\]
los coeficientes salen por que $dC_{I}\frac{\partial}{\partial x_{j}}=\frac{\partial}{\partial x_{j}}C_{I}$ y esto vale por \textasciicircum{}311f86 4. Por linealidad del diferencial de funciones suaves, claramente $d_U$ se extiende a una transformacion lineal sobre todo $\Omega^*(U)$ 5. Ademas $d_U$ satisface \textbf{(i)} por que si reemplazamos la expresion que dimos en 3. llegamos a

\[
d_U\omega:=\sum_{j=1}^{n}  \sum_{i_{1}< \ldots<i_{k}}\frac{\partial C_{i_1\ldots i_k}}{\partial x_j}\,dx_j\wedge dx_{i_1}\wedge\cdots\wedge dx_{i_k},
\]
y luego de permutar y posiblemente agregar un $-1$ nos queda que esto es una suma de $1$ hasta $n$ de $k+1$-formas 6. \textbf{(ii)} es trivial por que $f=f.1$ entonces $d|_{U}f=d|_{U}f.1=d|_{U}f$\\
7. \textbf{(iii)} Como $d_U$ es lineal y $\wedge$ es bilineal, basta con analizar cuando $\omega=f\,dx_{i_1}\wedge\cdots\wedge dx_{i_k}$ y $\theta=g\,dx_{j_1}\wedge\cdots\wedge dx_{j_\ell}$. Esto es por que el producto cuña de dos elementos de la suma directa se define como la suma de los productos cuña de los pares son $k$-formas del mismo grado $k$. 8. Luego por linealidad seria diferencial de cada uno de esos productos. Pero esos productos son basicamente como en multiplicacion de polinomios la suma de todas las combinaciones, y luego por linealidad nuevamente basta verlo para un caso cualquiera de esos que seria $\omega\wedge\theta$ 9. Entonces

\[
\begin{aligned}d_U(\omega\wedge\theta)& =d_U(fg\,dx_{i_1}\wedge\cdots\wedge dx_{i_k}\wedge dx_{j_1}\wedge\cdots\wedge dx_{j_\ell})\\ & =d(fg)\wedge dx_{i_1}\wedge\cdots\wedge dx_{i_k}\wedge dx_{j_1}\wedge\cdots\wedge dx_{j_\ell}\\&=(f\,dg+g\,df)\wedge dx_{i_1}\wedge\cdots\wedge dx_{i_k}\wedge dx_{j_1}\wedge\cdots\wedge dx_{j_\ell}\\& =f\,dg\wedge dx_{i_1}\wedge\cdots\wedge dx_{i_k}\wedge dx_{j_1}\wedge\cdots\wedge dx_{j_\ell}+g\,df\wedge dx_{i_1}\wedge\cdots\wedge dx_{i_k}\wedge dx_{j_1}\wedge\cdots\wedge dx_{j_\ell}\\&=(-1)^k f\,dx_{i_1}\wedge\cdots\wedge dx_{i_k}\wedge dg\wedge dx_{j_1}\wedge\cdots\wedge dx_{j_\ell}+df\wedge dx_{i_1}\wedge\cdots\wedge dx_{i_k}\wedge g\,dx_{j_1}\wedge\cdots\wedge dx_{j_\ell}\\&=(d_U\omega)\wedge\theta+(-1)^k\omega\wedge d_U\theta.\end{aligned}
\]
\begin{enumerate}[start=10]
\item \textbf{(iv)} $d_U^2=0$. Por linealidad de $d_U$, es suficiente con ver $d_Ud_U\omega=0$ con $\omega=f\,dx_I$, donde $dx_I=dx_{i_1}\wedge\cdots\wedge dx_{i_k}$.
\end{enumerate}
\[
\begin{aligned}d_Ud_U\omega&=d_U\left(\sum_j\frac{\partial f}{\partial x_j}dx_j\wedge dx_I\right)\\&=\sum_j d\left(\frac{\partial f}{\partial x_j}\right)\wedge dx_j\wedge dx_I\\&=\sum_j\sum_k\frac{\partial^2 f}{\partial x_k\partial x_j}dx_k\wedge dx_j\wedge dx_I\\&=\sum_{1\le j<k\le n}\left(\frac{\partial^2 f}{\partial x_k\partial x_j}-\frac{\partial^2 f}{\partial x_j\partial x_k}\right)dx_k\wedge dx_j\wedge dx_I=0.\end{aligned}
\]
Esto último vale porque

\[
\left(\frac{\partial^2 f}{\partial x_k\partial x_j}-\frac{\partial^2 f}{\partial x_j\partial x_k}\right)=\left[\frac{\partial }{\partial x_k},\frac{\partial}{\partial x_j}\right]f=0.
\]
\begin{enumerate}[start=11]
\item Por ultimo, definimos de manera global $d$. Sea $\theta\in\Omega^r(M)$. Dado $p\in M$, definimos
\end{enumerate}
\[
(d\theta)_p:=(d_U\theta|_U)_p
\]
para cualquier carta suave $(U,\varphi=(x_1,\ldots,x_n))$ de $M$ alrededor de $p$. (la parte. dela derecha cae en lo que hicimos al principio por lo tanto cumple todo) 12. Necesitamos ver que esta bien definido. Sea $(V,\psi=(y_1,\ldots,y_n))$ otra carta suave de $p$, y $W=U\cap V\ne\varnothing$. Por linealidad de la derivada exterior, alcanza con verlo cuando $\theta\in\Omega^k(M)$. 13. Tenemos

\[
\theta|_U=\sum_Ia_I\,dx_I,\qquad \theta|_V=\sum_Jb_J\,dy_J.
\]
\begin{enumerate}[start=14]
\item Entonces por definicion de $d_U$ y $d_V$,
\end{enumerate}
\[
d_U\theta|_U=\sum_Ida_I\wedge dx_I,\qquad d_V\theta|_V=\sum_Jdb_J\wedge dy_J.
\]
\begin{enumerate}[start=15]
\item Pero en $W=U\cap V$ hay una unica diferencial exterior local
\end{enumerate}
\[
d_W:\Omega^*(W)\to\Omega^{*}(W).
\]
\begin{enumerate}[start=16]
\item Tomando las dos expresiones de $\theta|_W$, esto ultimo es igual a, por ser $d_W$ diferencial exterior,
\end{enumerate}
\[
\sum_Ida_I\wedge dx_I=\sum_Jdb_J\wedge dy_J.
\]
\begin{enumerate}[start=17]
\item En particular, si evaluamos en $p$,
\end{enumerate}
\[
(\sum_Ida_I\wedge dx_I)_p=(\sum_Jdb_J\wedge dy_J)_p.
\]
\bookqed
\end{bookproof}
\section{Complejo de de Rham}
\begin{bookproposition}{Propiedades del pullback}
Sea $F:M\to N$ suave.

\begin{itemize}
\item (a) Si $\omega\in\Omega^k(N)$, entonces $F^*\omega\in\Omega^k(M)$, y $F^*:\Omega^k(N)\to\Omega^k(M)$ es transf. lineal.
\item (b) $F^*(\omega\wedge\theta)=(F^*\omega)\wedge(F^*\theta)$.
\item (c) $F^*(d_N\omega)=d_MF^*\omega$.
\end{itemize}
\end{bookproposition}
\begin{bookproof}{Ejercicio}
\begin{itemize}
\item $(a)$
\end{itemize}
\begin{enumerate}
\item Probemos que la sección $p\mapsto(F^*\omega)_p$ es suave. Fijemos $p\in M$ y tomemos cartas $(U,x)$ de $M$ alrededor de $p$ y $(V,y)$ de $N$ alrededor de $F(p)$, achicando $U$ de modo que $F(U)\subseteq V$. En $V$ escribimos
\end{enumerate}
\[
\omega=\sum_I a_I\,dy^{i_1}\wedge\cdots\wedge dy^{i_k},
\]
donde cada $a_I$ es suave. Afirmamos que sobre $U$ se cumple

\[
F^*\omega=\sum_I(a_I\circ F)\,d(y^{i_1}\circ F)\wedge\cdots\wedge d(y^{i_k}\circ F).
\]
\begin{enumerate}[start=2]
\item Para verificar esta igualdad, sean $q\in U$ y $v_1,\ldots,v_k\in T_qM$. Por definición del pullback,
\end{enumerate}
\[
\begin{aligned}(F^*\omega)_q(v_1,\ldots,v_k)&=\omega_{F(q)}\bigl((dF)_q(v_1),\ldots,(dF)_q(v_k)\bigr)\\&=\sum_I a_I(F(q))(dy^{i_1}\wedge\cdots\wedge dy^{i_k})_{F(q)}\bigl((dF)_q(v_1),\ldots,(dF)_q(v_k)\bigr).\end{aligned}
\]
\begin{enumerate}[start=3]
\item Por definición del producto wedge de $1$-formas,
\end{enumerate}
\[
(dy^{i_1}\wedge\cdots\wedge dy^{i_k})_{F(q)}\bigl((dF)_q(v_1),\ldots,(dF)_q(v_k)\bigr)=\det\left(dy^{i_r}_{F(q)}\bigl((dF)_q(v_s)\bigr)\right)_{r,s}.
\]
\begin{enumerate}[start=4]
\item Por la regla de la cadena, $dy^{i_r}_{F(q)}\circ(dF)_q=d(y^{i_r}\circ F)_q$. Por tanto,
\end{enumerate}
\[
\begin{aligned}\det\left(dy^{i_r}_{F(q)}\bigl((dF)_q(v_s)\bigr)\right)_{r,s}&=\det\left(d(y^{i_r}\circ F)_q(v_s)\right)_{r,s}\\&=\bigl(d(y^{i_1}\circ F)\wedge\cdots\wedge d(y^{i_k}\circ F)\bigr)_q(v_1,\ldots,v_k).\end{aligned}
\]
\begin{enumerate}[start=5]
\item Sustituyendo, obtenemos
\end{enumerate}
\[
\begin{aligned}(F^*\omega)_q(v_1,\ldots,v_k)&=\sum_I(a_I\circ F)(q)\bigl(d(y^{i_1}\circ F)\wedge\cdots\wedge d(y^{i_k}\circ F)\bigr)_q(v_1,\ldots,v_k)\\&=\left(\sum_I(a_I\circ F)\,d(y^{i_1}\circ F)\wedge\cdots\wedge d(y^{i_k}\circ F)\right)_q(v_1,\ldots,v_k).\end{aligned}
\]
\begin{enumerate}[start=6]
\item Como esto vale para todo $q\in U$ y todos $v_1,\ldots,v_k\in T_qM$, concluimos que
\end{enumerate}
\[
F^*\omega=\sum_I(a_I\circ F)\,d(y^{i_1}\circ F)\wedge\cdots\wedge d(y^{i_k}\circ F)
\]
sobre $U$.

\begin{enumerate}[start=7]
\item Cada función $a_I\circ F$ es suave y cada $d(y^{i_r}\circ F)$ es una $1$-forma suave. Por lo tanto, el lado derecho es una $k$-forma suave sobre $U$. Como $p$ era arbitrario, $F^*\omega\in\Omega^k(M)$.
\item Finalmente, sean $\omega,\eta\in\Omega^k(N)$ y $a,b\in\mathbb R$. Para todo $p\in M$ y $v_1,\ldots,v_k\in T_pM$,
\end{enumerate}
\[
\begin{aligned}\bigl(F^*(a\omega+b\eta)\bigr)_p(v_1,\ldots,v_k)&=(a\omega+b\eta)_{F(p)}\bigl((dF)_p(v_1),\ldots,(dF)_p(v_k)\bigr)\\&=a\,\omega_{F(p)}\bigl((dF)_p(v_1),\ldots,(dF)_p(v_k)\bigr)\\&\quad+b\,\eta_{F(p)}\bigl((dF)_p(v_1),\ldots,(dF)_p(v_k)\bigr)\\&=\bigl(aF^*\omega+bF^*\eta\bigr)_p(v_1,\ldots,v_k).\end{aligned}
\]
Luego $F^*(a\omega+b\eta)=aF^*\omega+bF^*\eta$, por lo que $F^*$ es lineal.

\begin{itemize}
\item $(b)$
\end{itemize}
\begin{enumerate}
\item Para probar que $F^*(\omega\wedge\theta)=(F^*\omega)\wedge(F^*\theta)$, evaluamos ambos lados en $p\in M$ y en vectores $v_1,\ldots,v_{k+\ell}\in T_pM$.
\item Por definición de pullback,
\end{enumerate}
\[
(F^*(\omega\wedge\theta))_p(v_1,\ldots,v_{k+\ell})=(\omega\wedge\theta)_{F(p)}((dF)_pv_1,\ldots,(dF)_pv_{k+\ell}).
\]
\begin{enumerate}[start=3]
\item Usando la definición de producto cuña,
\end{enumerate}
\[
(\omega\wedge\theta)_{F(p)}((dF)_pv_1,\ldots,(dF)_pv_{k+\ell})=\frac{1}{k!\ell!}\sum_{\pi\in S_{k+\ell}}\operatorname{sgn}(\pi)\,\omega_{F(p)}((dF)_pv_{\pi(1)},\ldots,(dF)_pv_{\pi(k)})\,\theta_{F(p)}((dF)_pv_{\pi(k+1)},\ldots,(dF)_pv_{\pi(k+\ell)}).
\]
\begin{enumerate}[start=4]
\item Pero, otra vez por definición de pullback,
\end{enumerate}
\[
\omega_{F(p)}((dF)_pv_{\pi(1)},\ldots,(dF)_pv_{\pi(k)})=(F^*\omega)_p(v_{\pi(1)},\ldots,v_{\pi(k)}),
\]
y

\[
\theta_{F(p)}((dF)_pv_{\pi(k+1)},\ldots,(dF)_pv_{\pi(k+\ell)})=(F^*\theta)_p(v_{\pi(k+1)},\ldots,v_{\pi(k+\ell)}).
\]
\begin{enumerate}[start=5]
\item Entonces
\end{enumerate}
\[
(F^*(\omega\wedge\theta))_p(v_1,\ldots,v_{k+\ell})=((F^*\omega)\wedge(F^*\theta))_p(v_1,\ldots,v_{k+\ell}).
\]
Como esto vale para todo $p$ y todo $v_1,\ldots,v_{k+\ell}$, queda

\[
F^*(\omega\wedge\theta)=(F^*\omega)\wedge(F^*\theta).
\]
\begin{itemize}
\item $(c)$
\end{itemize}
\begin{enumerate}
\item Por linealidad de $F^*$ y de la derivada exterior, es suficiente con verificar la igualdad en formas expresadas sobre un abierto coordenado $U$.
\item Recordemos que por definición de pullback de $0$-formas,
\item 
\end{enumerate}
\[
F^*(f)=f\circ F.
\]
\begin{enumerate}[start=4]
\item Así pues, dado $p\in M$ y $v\in T_pM$,
\end{enumerate}
\[
(F^*(df))_pv=(df)_{F(p)}((dF)_pv)=d(f\circ F)_pv.
\]
\begin{enumerate}[start=5]
\item Si $\theta$ es una $k$-forma, como la definición de la derivada exterior es local y no depende de la carta, fijemos un $p\in M$ y una carta $(U,\varphi=(x_1,\ldots,x_n))$ de $F(p)$:
\end{enumerate}
\[
\theta|_U=\sum_{i_1<\cdots<i_k}a_{i_1\ldots i_k}\,dx_{i_1}\wedge\cdots\wedge dx_{i_k}=\sum_Ia_I\,dx_I.
\]
\begin{enumerate}[start=6]
\item Entonces
\end{enumerate}
\[
(d_N\theta)|_U=\sum_I(da_I)\wedge dx_I.
\]
\begin{enumerate}[start=7]
\item Por tanto
\end{enumerate}
\[
\begin{aligned}(F^*(d_N\theta))_p&=\left(F^*\left(\sum_I(da_I)\wedge dx_I\right)\right)_p\\&=\sum_I(F^*(da_I))_p\wedge(F^*(dx_I))_p\\&=\sum_{i_1<\cdots<i_k}d(a_I\circ F)_p\wedge d(x_{i_1}\circ F)_p\wedge\cdots\wedge d(x_{i_k}\circ F)_p.\end{aligned}
\]
En la última igualdad usamos el paso 3.

\begin{enumerate}[start=8]
\item Por otro lado
\end{enumerate}
\[
\begin{aligned} d_M(F^*\theta))_p&=\left(d_M\left(F^*\left(\sum_Ia_I\,dx_I\right)\right)\right)_p\\&=\left(d_M\left(\sum_I(F^*a_I)\wedge F^*(dx_I)\right)\right)_p\\&=\left(d_M\left(\sum_{i_1<\cdots<i_k}(a_I\circ F)\,d(x_{i_1}\circ F)\wedge\cdots\wedge d(x_{i_k}\circ F)\right)\right)_p\\&=\sum_{i_1<\cdots<i_k}d(a_I\circ F)_p\wedge d(x_{i_1}\circ F)_p\wedge\cdots\wedge d(x_{i_k}\circ F)_p.\end{aligned}
\]
\begin{enumerate}[start=9]
\item En la última igualdad usamos que $d_Md_M=0$. Comparando con la expresión anterior, concluimos que
\end{enumerate}
\[
(F^*(d_N\theta))_p=(d_M(F^*\theta))_p.
\]
\bookqed
\end{bookproof}
\section{Distribuciones}
\begin{bookproposition}{D integrable implica involutiva}
Sea $\mathcal D$ una distribucion suave de dimension $k$. Si $\mathcal D$ es integrable, entonces es involutiva.

\end{bookproposition}
\begin{bookproof}{}
\begin{enumerate}
\item Sean $X,Y$ campos en $\mathcal D$. Queremos que $[X,Y]$ esta en $\mathcal D$.
\item Sea $q\in M$ arbitrario y sea $(N,F)$ subvariedad integral de $\mathcal D$ por $q=F(p)$, $p\in N$.
\item Como $(dF)_pT_pN=\mathcal D_{F(p)}$ para todo $p\in N$, y $X_{F(p)},Y_{F(p)}\in\mathcal D_{F(p)}$.
\item Por \hyperref[blk:b71da3]{\textasciicircum{}b71da3}, existen campos suaves $\widetilde X,\widetilde Y$ en $\mathfrak X(N)$ tales que
\end{enumerate}
\[
(dF)_p\widetilde X_p=X_{F(p)},\qquad (dF)_p\widetilde Y_p=Y_{F(p)},
\]
para todo $p\in N$. 5. Esto no es mas que $\widetilde X$ e $\widetilde Y$ estan $F$-relacionados con $X,Y$. Tenemos por teorema que $[\widetilde X,\widetilde Y]$ esta $F$-relacionado con $[X,Y]$; es decir

\[
(dF)_p[\widetilde X,\widetilde Y]_p=[X,Y]_{F(p)},
\]
para todo $p\in N$. 6. Por tanto, como $(dF)_p[\widetilde X,\widetilde Y]_p\in\mathcal D_{F(p)}$, por definicion de subvariedad integral, tenemos $[X,Y]_{F(p)}\in\mathcal D_{F(p)}$.

\bookqed
\end{bookproof}
\begin{booktheorem}{Frobenius local / Warner Th. 1.60}
Toda distribucion involutiva es integrable.

Mejor aun, para todo $p\in M$, existe una carta suave $(U,\varphi=(x_1,\ldots,x_n))$, cubico centrado en $p$, tal que las rebanadas

\[
D=\{q\in U:x_j(q)=\text{constante},\ j=k+1,\ldots,n\}
\]
son subvariedades integrales de $\mathcal D$.

Ademas, si $N\subset M$ es una subvariedad integral conexa de $\mathcal D$ tal que $N\subset U$, entonces $N$ debe estar contenida en una de estas rebanadas.

\end{booktheorem}
\begin{bookproof}{}
Sea $\mathcal D$ una distribucion de dimension $k$. Nos paramos en $p$, y tenemos $\mathcal D_p=\operatorname{span}\{v_1,\ldots,v_k\}$.

Sabemos por una proposicion probada en el capitulo de Campos que existe una carta $(V,\psi=(y_1,\ldots,y_n))$, centrada en $p$, tal que

\[
v_j=\frac{\partial}{\partial y_j}\bigg|_p,\qquad j=1,\ldots,k.
\]
Nos llevamos la distribucion para el abierto $\psi(V)$ de $\mathbb R^n$, usando la carta. Tenemos una distribucion $\mathcal D$ en $\psi(V)$, con $0=\psi(p)$, y tal que

\[
\mathcal D_0=\operatorname{span}\left\{\frac{\partial}{\partial y_1}\bigg|_0,\ldots,\frac{\partial}{\partial y_k}\bigg|_0\right\}.
\]
Como la distribucion es suave, tenemos campos suaves $X_1,\ldots,X_k$ definidos en un abierto de $\psi(V)$ tal que

\[
\mathcal D_q=\operatorname{span}\{(X_1)_q,\ldots,(X_k)_q\}.
\]
Por tanto, podemos usar el marco $\left\{\frac{\partial}{\partial y_1},\ldots,\frac{\partial}{\partial y_n}\right\}$ para escribir estos campos como combinacion lineal

\[
X_j=\sum_{i=1}^na_{ij}\frac{\partial}{\partial y_i},
\]
donde $a_{ij}$ son funciones suaves sobre un abierto contenido en $\psi(V)$.

Armando la matriz $T(q)=(a_{ij}(q))$, cuando hacemos $T(0)$, esta matriz tiene la submatriz $T(0)_{1\ldots k,1\ldots k}$ igual a la identidad, porque $(X_j)_0=\frac{\partial}{\partial y_j}\big|_0$. Podemos pensar en la funcion

\[
f:\psi(V)\to\mathbb R,\qquad q\mapsto\det(T(q)_{1\ldots k,1\ldots k}).
\]
La cual es una funcion suave por las entradas de la matriz $T(q)$ son suaves. Como $f(0)=1\ne0$, existe un abierto de $0$ donde $f(q)\ne0$. Achicamos $\psi(V)$, si hace falta, y podemos asumir que esto pasa en $\psi(V)$.

Por tanto, tiene sentido hablar de la matriz inversa de la submatriz $T(q)_{1\ldots k,1\ldots k}$, digamos $Q(q)$. Esto define una funcion

\[
Q:\psi(V)\to M_{k\times k}(\mathbb R),\qquad q\mapsto Q(q),
\]
tal que

\[
T(q)_{1\ldots k,1\ldots k}\cdot Q(q)=I_k.
\]
Notar que las entradas de la matriz que da $Q$ son funciones suaves, pues son funciones racionales entre productos y sumas con las entradas de $T(q)$, con denominador el determinante de $T(q)_{1\ldots k,1\ldots k}$, el cual no se anula en $\psi(V)$.

Luego

\[
T(q)Q(q)=\begin{pmatrix}1& &0\\ &\ddots& \\0& &1\\ b_{k+1,1}&\cdots&b_{k+1,k}\\ \vdots& &\vdots\\ b_{n1}&\cdots&b_{nk}\end{pmatrix}.
\]
¿Que hemos hecho? Por algebra, lo que se ha hecho es cambiar los generadores de la distribucion. Como los vectores columnas de $T(q)$ generan $\mathcal D_q$ y $Q(q)$ da los escalares para que las combinaciones lineales de las columnas de $T(q)$ generen los vectores columnas de la matriz de la derecha, tenemos que $\mathcal D$ es generada en el abierto $\psi(V)$ por los campos $\widetilde X_1,\ldots,\widetilde X_k$ con

\[
\widetilde X_j=\frac{\partial}{\partial y_j}+\sum_{i=k+1}^nb_{ij}\frac{\partial}{\partial y_i},\qquad b_{ij}\in C^\infty(V).
\]
La forma de pensarlo: recordemos que en algebra 2, cuando pedian dar una base de un subespacio vectorial generado por $k$ vectores, se ponian en una matriz como vectores fila y se hacia Gauss-Jordan; en cada paso estamos haciendo combinaciones lineales de las filas, y el resultado final es otro conjunto generador muy lindo del subespacio vectorial.

Lo lindo que tiene este nuevo conjunto de generadores es que conmutan.

Hagamos:

\[
[\widetilde X_i,\widetilde X_j]=\left[\frac{\partial}{\partial y_i}+\sum_{h=k+1}^nb_{hi}\frac{\partial}{\partial y_h},\frac{\partial}{\partial y_j}+\sum_{\ell=k+1}^nb_{\ell j}\frac{\partial}{\partial y_\ell}\right].
\]
Por bilinealidad del corchete y que $\left[\frac{\partial}{\partial y_i},\frac{\partial}{\partial y_j}\right]=0$, queda

\[
[\widetilde X_i,\widetilde X_j]=\left[\frac{\partial}{\partial y_i},\sum_{\ell=k+1}^nb_{\ell j}\frac{\partial}{\partial y_\ell}\right]-\left[\frac{\partial}{\partial y_j},\sum_{h=k+1}^nb_{hi}\frac{\partial}{\partial y_h}\right]+\left[\sum_{h=k+1}^nb_{hi}\frac{\partial}{\partial y_h},\sum_{\ell=k+1}^nb_{\ell j}\frac{\partial}{\partial y_\ell}\right].
\]
Esto pertenece a

\[
\operatorname{span}\left\{\frac{\partial}{\partial y_{k+1}},\ldots,\frac{\partial}{\partial y_n}\right\}.
\]
Por hipotesis, la distribucion es involutiva, por tanto $[\widetilde X_i,\widetilde X_j]\in\mathcal D$. Y como

\[
\mathcal D_q\cap\operatorname{span}\left\{\frac{\partial}{\partial y_{k+1}}\bigg|_q,\ldots,\frac{\partial}{\partial y_n}\bigg|_q\right\}=\{0\},
\]
se sigue la afirmacion

\[
[\widetilde X_i,\widetilde X_j]=0.
\]
Con estos campos vamos a construir de manera natural una subvariedad integral. ¿Por que natural? Como los campos estan en la distribucion, entonces sus curvas integrales van a estar contenidas en las subvariedades integrales. Lo natural es combinar los flujos de cada campo para ir formando la subvariedad. Sabemos que esto sale bien cuando los campos conmutan.

Si en estos momentos tuvieramos que dar una subvariedad integral por el origen $0$, usando que podemos tomar un abierto cubico de $0$ donde hay uniformidad conjunta para $\widetilde X_1,\ldots,\widetilde X_k$, con flujos $\theta^1,\ldots,\theta^k$ respectivamente, definimos

\[
f:(-\delta,\delta)^k\to\mathbb R^n,\qquad (t_1,\ldots,t_k)\mapsto\theta^1_{t_1}\circ\cdots\circ\theta^k_{t_k}(0).
\]
\bookqed
\end{bookproof}
\section{Orientacion de espacios vectoriales}
\begin{bookproposition}{Caracterizacion de orientabilidad}
Sean $M$ una variedad suave conexa de dimension $n$. Entonces las siguientes afirmaciones son equivalentes:

\begin{itemize}
\item (a) $M$ es orientable.
\item (b) Existe una $n$-forma continua nunca nula sobre $M$.
\item (c) Existe una $n$-forma suave nunca nula sobre $M$.
\end{itemize}
\end{bookproposition}
\begin{bookproof}{}
\begin{itemize}
\item \textbf{$(c)\Rightarrow(b)$} es inmediato.
\item \textbf{$(b)\Rightarrow(a)$}.
\end{itemize}
\begin{enumerate}
\item Tenemos $\omega$ una $n$-forma no nula, recordamos que un $n$-tensor alternante no nulo define una orientación (en espacios vectoriales)
\item Usamos la continuidad de la $n$-forma $\omega$ para definir un atlas de orientación.
\item Arrancamos con $\mathcal A=\{(U_\alpha,\varphi_\alpha=(x_1^\alpha,\ldots,x_n^\alpha))\}_{\alpha\in I}$ un atlas suave de $M$, donde podemos suponer que cada $U_\alpha$ es un conjunto conexo.
\item Recordar que las componentes conexas en un espacio topológico localmente conexo por caminos son abiertas.
\item Para cada $(U_\alpha,\varphi_\alpha)\in\mathcal A$, consideramos la función
\end{enumerate}
\[
f_\alpha:U_\alpha\to\mathbb R\qquad p\mapsto \omega_p\!\left(\left.\frac{\partial}{\partial x_1^\alpha}\right|_p,\ldots,\left.\frac{\partial}{\partial x_n^\alpha}\right|_p\right)
\]
\begin{enumerate}[start=6]
\item Como $\omega$ es nunca nula y continua, entonces $f_\alpha$ es continua y nunca nula.
\item Como $U_\alpha$ es conexo, se sigue que $f_\alpha>0$ o bien $f_\alpha<0$ en todo $U_\alpha$.
\item Si $f_\alpha<0$, construimos la nueva carta suave $(U_\alpha,\widetilde\varphi_\alpha=(-x_1^\alpha,x_2^\alpha,\ldots,x_n^\alpha))$.
\item Nuestro atlas de orientación es
\end{enumerate}
\[
\mathcal A_0=\{(U_\alpha,\varphi_\alpha)\in\mathcal A:f_\alpha>0\}\cup\{(U_\alpha,\widetilde\varphi_\alpha):(U_\alpha,\varphi_\alpha)\in\mathcal A,\ f_\alpha<0\}.
\]
\begin{itemize}
\item \textbf{$\mathcal A_0$ es un atlas suave de $M$}: Sus cartas son suavemente compatibles. Veamoslo:
\end{itemize}
\begin{enumerate}
\item Si $(U_\alpha,\varphi_\alpha),(U_\beta,\varphi_\beta)$ son t.q. $f_\alpha>0$, $f_\beta>0$ y $U_\alpha\cap U_\beta\neq\varnothing$, entonces son suavemente compatibles porque están en $\mathcal A$. Y $\mathcal A$ es atlas suave
\item Si $(U_\alpha,\varphi_\alpha)$ y $(U_\beta,\widetilde\varphi_\beta)$ donde $f_\alpha>0$, $f_\beta<0$ y $U_\alpha\cap U_\beta\neq\varnothing$. Consideremos
\end{enumerate}
\[
T:\mathbb R^n\to\mathbb R^n,\qquad (x^1,\ldots,x^n)\mapsto(-x^1,x^2,\ldots,x^n),
\]
la cual es una transf. lineal t.q. $T^2=\operatorname{Id}_{\mathbb R^n}$ y

\[
\widetilde\varphi_\beta=T\circ\varphi_\beta,
\]
así

\[
\widetilde\varphi_\beta\circ\varphi_\alpha^{-1}=T\circ \underbrace{\varphi_\beta\circ\varphi_\alpha^{-1}}_{\text{suave, pues }\mathcal A\text{ es atlas}}
\]
y

\[
\varphi_\alpha\circ\widetilde\varphi_\beta^{-1}=\varphi_\alpha\circ(T\circ\varphi_\beta)^{-1}=\underbrace{\varphi_\alpha\circ\varphi_\beta^{-1}}_{\text{suave, pues }\mathcal A\text{ es atlas}}\circ T^{-1},
\]
\begin{enumerate}[start=3]
\item Y finalmente si $(U_\alpha,\widetilde\varphi_\alpha)$, $(U_\beta,\widetilde\varphi_\beta)$ son t.q $f_{\alpha }<0$ y $f_{\beta}<0$ entonces son suavemente compatibles:
\end{enumerate}
\[
\widetilde\varphi_\beta\circ\widetilde\varphi_\alpha^{-1}=(T\circ\varphi_\beta)\circ(\varphi_\alpha^{-1}\circ T^{-1})=T\circ(\varphi_\beta\circ\varphi_\alpha^{-1})\circ T^{-1}
\]
\begin{itemize}
\item \textbf{Para todo $p\in M$, las cartas de $\mathcal A_0$ definen la misma orientación en $T_pM$:}
\end{itemize}
\begin{enumerate}
\item Sean $(U,\varphi)$, $(V,\psi)\in\mathcal A_0$ t.q. $U\cap V\neq\varnothing$. Tenemos tres casos
\end{enumerate}
\begin{itemize}
\item \textbf{1er Caso:}

\begin{enumerate}
\item Primero tenemos
\end{enumerate}
\end{itemize}
\[
(U,\varphi)=(U_\alpha,\varphi_\alpha),\qquad (V,\psi)=(U_\beta,\varphi_\beta),\qquad f_\alpha>0,\ f_\beta>0.
\]
entonces

\[
0<f_\alpha=\omega\!\left(\frac{\partial}{\partial x_\alpha^1},\ldots,\frac{\partial}{\partial x_\alpha^n}\right)=\det\!\left(\frac{\partial x_\beta^j}{\partial x_\alpha^i}\right)\omega\!\left(\frac{\partial}{\partial x_\beta^1},\ldots,\frac{\partial}{\partial x_\beta^n}\right)=\det\!\left(\frac{\partial x_\beta^j}{\partial x_\alpha^i}\right)f_\beta.
\]
\begin{Verbatim}[fontsize=\small]
 2. Luego 
\end{Verbatim}
\[
\det\!\left(\frac{\partial x_\beta^j}{\partial x_\alpha^i}\right)>0.
\]
\begin{itemize}
\item \textbf{2do caso:}

\begin{enumerate}
\item Si $(U,\varphi)=(U_\alpha,\varphi_\alpha)$ y $(V,\psi)=(U_\beta,\widetilde\varphi_\beta)$, con $f_\alpha>0$ y $f_\beta<0$
\item Entonces, repitiendo la misma cuenta de arriba, tenemos
\end{enumerate}
\end{itemize}
\[
\begin{aligned} 0<f_\alpha & =\omega\!\left(\frac{\partial}{\partial x_\alpha^1},\ldots,\frac{\partial}{\partial x_\alpha^n}\right)\\ &=\det\!\left(\frac{\partial \widetilde x_\beta^j}{\partial x_\alpha^i}\right)\omega\!\left(\frac{\partial}{\partial \widetilde x_\beta^1},\ldots,\frac{\partial}{\partial \widetilde x_\beta^n}\right)\\&=\det\!\left(\frac{\partial \widetilde x_\beta^j}{\partial x_\alpha^i}\right)\omega\!\left(-\frac{\partial}{\partial x_\beta^1},\ldots,\frac{\partial}{\partial x_\beta^n}\right)\\&=-\det\!\left(\frac{\partial \widetilde x_\beta^j}{\partial x_\alpha^i}\right)\omega\!\left(\frac{\partial}{\partial x_\beta^1},\ldots,\frac{\partial}{\partial x_\beta^n}\right)\\&=-\det\!\left(\frac{\partial \widetilde x_\beta^j}{\partial x_\alpha^i}\right)f_{\beta }\end{aligned}
\]
\begin{Verbatim}[fontsize=\small]
 3. Pero como $0>f_{\beta}$ tenemos que 
\end{Verbatim}
\[
\det\!\left(\frac{\partial \widetilde x_\beta^j}{\partial x_\alpha^i}\right)>0
\]
\begin{itemize}
\item \textbf{3er Caso}

\begin{enumerate}
\item Por último, si $(U,\varphi)=(U_\alpha,\widetilde\varphi_\alpha)$ y $(V,\psi)=(U_\beta,\widetilde\varphi_\beta)$ con $f_\alpha<0$ y $f_\beta<0$,
\item Entonces tenemos
\end{enumerate}
\end{itemize}
\[
\begin{aligned} 0>f_\alpha & =\omega\!\left(\frac{\partial}{\partial x_\alpha^1},\ldots,\frac{\partial}{\partial x_\alpha^n}\right)\\&=-\omega\!\left(\frac{\partial}{\partial \widetilde x_\alpha^1},\ldots,\frac{\partial}{\partial \widetilde x_\alpha^n}\right)\\&=\det\!\left(\frac{\partial \widetilde x_\beta^j}{\partial \widetilde x_\alpha^i}\right)-\omega\!\left(\frac{\partial}{\partial \widetilde x_\beta^1},\ldots,\frac{\partial}{\partial \widetilde x_\beta^n}\right)\\& =\det\!\left(\frac{\partial \widetilde x_\beta^j}{\partial \widetilde x_\alpha^i}\right)\omega\!\left(\frac{\partial}{\partial  x_\beta^1},\ldots,\frac{\partial}{\partial  x_\beta^n}\right)\\&=\det\!\left(\frac{\partial \widetilde x_\beta^j}{\partial \widetilde x_\alpha^i}\right)f_{\beta }\end{aligned}
\]
\begin{Verbatim}[fontsize=\small]
 3. Con lo cual 
\end{Verbatim}
\[
\det\!\left(\frac{\partial \widetilde x_\beta^j}{\partial \widetilde x_\alpha^i}\right)>0
\]
\begin{enumerate}[start=4]
\item Luego todas las funciones de transición de $\mathcal A_0$ tienen determinante positivo. Por consiguiente, todas las cartas de $\mathcal A_0$ determinan la misma orientación.
\end{enumerate}
\begin{itemize}
\item \textbf{$(a)\Rightarrow (c)$}
\end{itemize}
\begin{enumerate}
\item Tenemos el atlas de orientación $\mathcal A$. Queremos construir $\omega\in\Omega^n(M)$.
\item Para cada $(U_\alpha,\varphi_\alpha)\in\mathcal A$, con $\varphi_\alpha=(x_\alpha^1,\ldots,x_\alpha^n)$, consideremos
\end{enumerate}
\[
\xi_\alpha=(dx_\alpha^1)\wedge\cdots\wedge(dx_\alpha^n).
\]
\begin{enumerate}[start=3]
\item Tomemos una partición de la unidad subordinada al cubrimiento $\{U_\alpha\}_{\alpha\in I}$, digamos $\{\rho_\alpha\}_{\alpha\in I}$.
\item Definimos sobre $M$
\end{enumerate}
\[
\rho_\alpha\,\xi_\alpha=\begin{cases}\rho_\alpha(q)\,\xi_\alpha(q),&q\in U_\alpha,\\0,&q\notin U_\alpha.\end{cases}
\]
\begin{enumerate}[start=5]
\item Como $\rho_\alpha$ es suave, usando el criterio de suavidad para formas, tenemos que $\rho_\alpha\,\xi_\alpha$ es suave.
\item Definimos entonces
\end{enumerate}
\[
\omega=\sum_{\alpha\in I}\rho_\alpha\,\xi_\alpha.
\]
\begin{enumerate}[start=7]
\item Y $\omega$ resulta suave, pues $\{\operatorname{supp}\rho_\alpha\}_{\alpha\in I}$ es un conjunto localmente finito luego para cada $p\in M$ existe un entorno abierto $V$ de $p$ tal que sólo un número finito de las formas $\rho_\alpha\xi_\alpha$ son no nulas en $V$. En consecuencia, $\omega|_V$ es una suma finita de formas suaves, y por tanto es suave.
\item Además, para cada $p\in M$, $\omega_{p}$ es $n$-lineal y alternante.
\item Sólo resta verificar que $\omega$ nunca se anula. Para esto recordemos que $\sum_{\alpha\in I}\rho_\alpha(p)=1$ y $\rho_\alpha(p)\ge0$ para todo $\alpha\in I$.
\item Dado $p\in M$ como $\sum_{\alpha\in I}\rho_\alpha(p)=1$ existe $\beta\in I$ tal que $\rho_\beta(p)>0$ osea $p\in \operatorname{supp} \rho_{\alpha }\subseteq U_{\beta}$
\item La carta $(U_\beta,\varphi_\beta=(x_\beta^1,\ldots,x_\beta^n))$ determina la base $\{\frac{\partial}{\partial x_\beta^1}\big|_p,\ldots,\frac{\partial}{\partial x_\beta^n}\big|_p\}$ de $T_pM$.
\item Entonces
\end{enumerate}
\[
\begin{aligned}\omega_{p}\!\left(\frac{\partial}{\partial x_\beta^1}\bigg|_p,\ldots,\frac{\partial}{\partial x_\beta^n}\bigg|_p\right)&= \rho_{\beta }(p)\xi_{\beta}\!\left(\frac{\partial}{\partial x_\beta^1}\bigg|_p,\ldots,\frac{\partial}{\partial x_\beta^n}\bigg|_p\right)+\sum_{\substack{\alpha\in I \\ \alpha  \neq \beta }}^{n} \rho_\alpha(p)\,\xi_{\alpha}\!\left(\frac{\partial}{\partial x_\beta^1}\bigg|_p,\ldots,\frac{\partial}{\partial x_\beta^n}\bigg|_p\right)\\&=\rho_{\beta }(p)+\sum_{\substack{\alpha\in I \\ \alpha  \neq \beta }}^{n} \rho_\alpha(p)\det\left(\frac{\partial x_{\alpha }^{j } }{\partial x_{\beta }^{i}}\right) \end{aligned}
\]
para el segundo igual usamos la definicion de evaluar una $n$-forma y la propiedad $dx(\frac{\partial }{\partial y})=\frac{\partial }{\partial y}x$ 13. Y el resultado de esta suma es no negativo dado que $\rho_{\beta}(p)>0$ y $\rho_{\alpha }\geq 0$ y el $\det$ es mayor que $0$ por que ambas cartas pertencen a $\mathcal A$ el atlas de orientacion

\bookqed
\end{bookproof}
